\documentclass[12pt,aspectratio=1610]{beamer}
\geometry{paperwidth=20cm,paperheight=12.5cm}
\usepackage[italian,provide=*]{babel}
\usepackage{lmodern,fontspec}
\setsansfont{NimbusSans-Regular.otf}[Path=sorgenti/font/,BoldFont=NimbusSans-Bold.otf,ItalicFont=NimbusSans-Italic.otf]
\setmainfont{NimbusSans-Regular.otf}[Path=sorgenti/font/,BoldFont=NimbusSans-Bold.otf,ItalicFont=NimbusSans-Italic.otf]
\setmonofont{lmmono10-regular.otf}[Path=sorgenti/font/,Scale=.88]
\usefonttheme{professionalfonts}
\usepackage{amsmath,amssymb,booktabs,array,tabularx,graphicx,tikz,listings,bookmark}
\usetikzlibrary{calc,arrows.meta,positioning}
\definecolor{Paper}{HTML}{F8F6F1}
\definecolor{Ink}{HTML}{253344}\definecolor{Navy}{HTML}{15283E}
\definecolor{Quantum}{HTML}{5146A5}\definecolor{Classical}{HTML}{087C74}
\definecolor{Measure}{HTML}{A05B10}\definecolor{Rose}{HTML}{9B4564}
\definecolor{Muted}{HTML}{657180}\definecolor{Rule}{HTML}{D6DAD9}
\definecolor{Wash}{HTML}{EAE7F6}\definecolor{Night}{HTML}{112236}
\definecolor{Lilac}{HTML}{ABB5FF}\definecolor{Mint}{HTML}{72D8C8}\definecolor{Gold}{HTML}{EEC184}
\setbeamercolor{background canvas}{bg=Paper}\setbeamercolor{normal text}{fg=Ink}
\setbeamertemplate{navigation symbols}{}\setbeamertemplate{headline}{}\setbeamertemplate{footline}{}
\setbeamersize{text margin left=0pt,text margin right=0pt}\setbeameroption{hide notes}
\newsavebox{\BodyBox}\newsavebox{\SideBox}\newsavebox{\RecallBox}\newsavebox{\TitleBox}
\newcommand{\MeasureBox}[2]{\typeout{GEOBOX:#1:\the\wd#2:\the\ht#2:\the\dp#2}}
\lstset{language=R,basicstyle=\ttfamily\fontsize{10.2}{13.4}\selectfont,
columns=fullflexible,keepspaces=true,showstringspaces=false,breaklines=true,breakatwhitespace=false,
commentstyle=\color{Muted},keywordstyle=\color{Classical},stringstyle=\color{Quantum},
aboveskip=0pt,belowskip=0pt,xleftmargin=0pt,xrightmargin=0pt,tabsize=2}
\hypersetup{colorlinks=true,linkcolor=Classical,urlcolor=Classical,pdfpagemode=UseOutlines,pdfstartview=Fit,bookmarksopen=true,bookmarksopenlevel=0}
\bookmarksetup{open,openlevel=0,depth=2}
\tikzset{arr/.style={-{Stealth[length=2mm]},line width=.8pt,draw=Classical}}
\hypersetup{pdftitle={Campionamento, stima e incertezza},pdfauthor={},pdfsubject={Corso di statistica e data science per Geologia, primo anno; modulo 05}}
\usetikzlibrary{decorations.pathreplacing,fit,backgrounds,matrix}\hypersetup{pdftitle={M13 · Cronologie e archivi sedimentari},pdfsubject={Revisione integrale R6}}\begin{document}
\begin{frame}[plain]\hypertarget{slide1}{}\begin{tikzpicture}[remember picture,overlay]\fill[Night] (current page.south west) rectangle (current page.north east);\draw[Mint,line width=2pt] ([xshift=1.05cm,yshift=-1.35cm]current page.north west)--([xshift=3cm,yshift=-1.35cm]current page.north west);\node[anchor=north west,inner sep=0] at ([xshift=1.05cm,yshift=-1.7cm]current page.north west){\color{Mint}\bfseries\fontsize{11}{14}\selectfont M13 · REVISIONE R6};
\node[anchor=north west,inner sep=0] at ([xshift=1.05cm,yshift=-2.65cm]current page.north west){\begin{minipage}{15.8cm}\raggedright\color{white}\bfseries\fontsize{31}{38}\selectfont Cronologie e archivi sedimentari\end{minipage}};
\node[anchor=north west,inner sep=0] at ([xshift=1.05cm,yshift=-6.4cm]current page.north west){\begin{minipage}{15.8cm}\raggedright\color{Lilac}\fontsize{15}{21}\selectfont Modelli età–profondità, proxy e propagazione dell’incertezza\end{minipage}};
\node[anchor=north west,inner sep=0] at ([xshift=1.05cm,yshift=-8.6cm]current page.north west){\begin{minipage}{16.5cm}\color{white}\fontsize{11}{15}\selectfont 30 slide · teoria, derivazioni, figure e applicazioni in R.\par\smallskip Provenienza dei dati e sorgenti nei materiali allegati.\end{minipage}};
\node[anchor=south east,inner sep=0] at ([xshift=18.95cm,yshift=-11.35cm]current page.north west){\color{Lilac!45}\bfseries\fontsize{64}{68}\selectfont M13};
\end{tikzpicture}\end{frame}
\begin{frame}[plain]\hypertarget{slide2}{}\bookmark[level=0,page=2]{Indice}\begin{tikzpicture}[remember picture,overlay]\node[anchor=north west,inner sep=0] at ([xshift=0.95cm,yshift=-0.65cm]current page.north west){\color{Navy}\bfseries\fontsize{22}{27}\selectfont Percorso del modulo};
\draw[Classical,line width=1.3pt] ([xshift=.95cm,yshift=-1.68cm]current page.north west)--([xshift=2.25cm,yshift=-1.68cm]current page.north west);\node[anchor=north west,inner sep=0] at ([xshift=0.95cm,yshift=-2.23cm]current page.north west){\color{Classical}\bfseries\fontsize{13}{16}\selectfont 3};
\node[anchor=north west,inner sep=0] at ([xshift=2.0cm,yshift=-2.23cm]current page.north west){\begin{minipage}{7.25cm}\raggedright\color{Navy}\bfseries\fontsize{12.4}{15.5}\selectfont\hyperlink{slide3}{Archivio, ancore e convenzioni temporali}\end{minipage}};
\node[anchor=north west,inner sep=0] at ([xshift=0.95cm,yshift=-3.5300000000000002cm]current page.north west){\color{Classical}\bfseries\fontsize{13}{16}\selectfont 7};
\node[anchor=north west,inner sep=0] at ([xshift=2.0cm,yshift=-3.5300000000000002cm]current page.north west){\begin{minipage}{7.25cm}\raggedright\color{Navy}\bfseries\fontsize{12.4}{15.5}\selectfont\hyperlink{slide7}{Interpolazione e tassi di accumulo}\end{minipage}};
\node[anchor=north west,inner sep=0] at ([xshift=0.95cm,yshift=-4.83cm]current page.north west){\color{Classical}\bfseries\fontsize{13}{16}\selectfont 11};
\node[anchor=north west,inner sep=0] at ([xshift=2.0cm,yshift=-4.83cm]current page.north west){\begin{minipage}{7.25cm}\raggedright\color{Navy}\bfseries\fontsize{12.4}{15.5}\selectfont\hyperlink{slide11}{Scenari cronologici e propagazione}\end{minipage}};
\node[anchor=north west,inner sep=0] at ([xshift=0.95cm,yshift=-6.130000000000001cm]current page.north west){\color{Classical}\bfseries\fontsize{13}{16}\selectfont 16};
\node[anchor=north west,inner sep=0] at ([xshift=2.0cm,yshift=-6.130000000000001cm]current page.north west){\begin{minipage}{7.25cm}\raggedright\color{Navy}\bfseries\fontsize{12.4}{15.5}\selectfont\hyperlink{slide16}{Proxy e risoluzione temporale}\end{minipage}};
\node[anchor=north west,inner sep=0] at ([xshift=10.299999999999999cm,yshift=-2.23cm]current page.north west){\color{Classical}\bfseries\fontsize{13}{16}\selectfont 21};
\node[anchor=north west,inner sep=0] at ([xshift=11.35cm,yshift=-2.23cm]current page.north west){\begin{minipage}{7.25cm}\raggedright\color{Navy}\bfseries\fontsize{12.4}{15.5}\selectfont\hyperlink{slide21}{Discontinuità e sensibilità del modello}\end{minipage}};
\node[anchor=north west,inner sep=0] at ([xshift=10.299999999999999cm,yshift=-3.5300000000000002cm]current page.north west){\color{Classical}\bfseries\fontsize{13}{16}\selectfont 25};
\node[anchor=north west,inner sep=0] at ([xshift=11.35cm,yshift=-3.5300000000000002cm]current page.north west){\begin{minipage}{7.25cm}\raggedright\color{Navy}\bfseries\fontsize{12.4}{15.5}\selectfont\hyperlink{slide25}{Pendenze, durate ed eventi}\end{minipage}};
\node[anchor=north west,inner sep=0] at ([xshift=10.299999999999999cm,yshift=-4.83cm]current page.north west){\color{Classical}\bfseries\fontsize{13}{16}\selectfont 30};
\node[anchor=north west,inner sep=0] at ([xshift=11.35cm,yshift=-4.83cm]current page.north west){\begin{minipage}{7.25cm}\raggedright\color{Navy}\bfseries\fontsize{12.4}{15.5}\selectfont\hyperlink{slide30}{Consegna e studio autonomo}\end{minipage}};
\node[anchor=north west,inner sep=0] at ([xshift=0.95cm,yshift=-10.55cm]current page.north west){\begin{minipage}{18cm}\color{Muted}\fontsize{10}{13}\selectfont I numeri indicano le pagine iniziali; i titoli sono cliccabili. Codice, figure e verifiche sono inclusi nel progetto.\end{minipage}};
\end{tikzpicture}\end{frame}
\begin{lrbox}{\TitleBox}\begin{minipage}[t]{18.2cm}\raggedright\sloppy\color{Ink}\fontsize{20.7}{24}\selectfont \bfseries\color{Navy}La cronologia collega profondità e tempo\end{minipage}\end{lrbox}
\begin{lrbox}{\BodyBox}\begin{minipage}[t]{10.25cm}\raggedright\sloppy\color{Ink}\fontsize{12.4}{16.4}\selectfont Un archivio sedimentario conserva osservazioni ordinate nello spazio stratigrafico, mentre l’interpretazione temporale richiede una relazione fra profondità ed età. Il capitolo utilizza sei ancore sintetiche distribuite fra zero e 160 centimetri, con età di calendario comprese fra zero e 2200 anni BP. Un proxy simulato è disponibile ogni due centimetri. Le deviazioni standard associate alle ancore definiscono scenari didattici di sensibilità, dei quali occorre conservare la natura ipotetica quando si interpretano bande e durate derivate.\par\vspace{.28cm}{\centering\fontsize{13}{17}\selectfont\(\displaystyle\begin{gathered}t=t(z),\qquad z\in[0{,}160]\ \mathrm{cm}\end{gathered}\)\par}\end{minipage}\end{lrbox}
\begin{lrbox}{\SideBox}\begin{minipage}[t]{7.45cm}\raggedright\sloppy\color{Ink}\fontsize{12.2}{16.1}\selectfont \includegraphics[width=7.45cm]{figure/M13/F01.pdf}\par\vspace{.16cm}{\fontsize{9.5}{12.6}\selectfont Sei ancore sintetiche e interpolazione lineare a tratti. La profondità aumenta verso il basso.}\end{minipage}\end{lrbox}
\begin{lrbox}{\RecallBox}\begin{minipage}[t]{16.45cm}\raggedright\sloppy\color{Ink}\fontsize{10.1}{13.2}\selectfont Il passaggio dalla profondità al tempo richiede un modello cronologico esplicito.\end{minipage}\end{lrbox}
\MeasureBox{3-TitleBox}{\TitleBox}
\MeasureBox{3-BodyBox}{\BodyBox}
\MeasureBox{3-SideBox}{\SideBox}
\MeasureBox{3-RecallBox}{\RecallBox}
\begin{frame}[plain]\hypertarget{slide3}{}
\bookmark[level=0,page=3]{Archivio, ancore e convenzioni temporali}
\bookmark[level=1,page=3]{La cronologia collega profondità e tempo}
\begin{tikzpicture}[remember picture,overlay]\node[anchor=north east,inner sep=0] at ([xshift=19.05cm,yshift=-0.16cm]current page.north west){\color{Classical}\fontsize{8}{10}\selectfont\hyperlink{slide2}{Indice}};
\node[anchor=north west,inner sep=0] at ([xshift=0.95cm,yshift=-0.62cm]current page.north west){\usebox{\TitleBox}};
\draw[Classical,line width=1.3pt] ([xshift=.95cm,yshift=-1.68cm]current page.north west)--([xshift=2.25cm,yshift=-1.68cm]current page.north west);\node[anchor=north west,inner sep=0] at ([xshift=0.95cm,yshift=-2.12cm]current page.north west){\usebox{\BodyBox}};
\node[anchor=north west,inner sep=0] at ([xshift=11.6cm,yshift=-2.1cm]current page.north west){\color{Quantum}\bfseries\fontsize{11.5}{14}\selectfont Evidenza grafica};
\node[anchor=north west,inner sep=0] at ([xshift=11.6cm,yshift=-2.75cm]current page.north west){\usebox{\SideBox}};
\draw[Rule,line width=.5pt] ([xshift=.95cm,yshift=-10.43cm]current page.north west)--([xshift=19.05cm,yshift=-10.43cm]current page.north west);\node[anchor=north west,inner sep=0] at ([xshift=0.95cm,yshift=-10.66cm]current page.north west){\color{Classical}\bfseries\fontsize{9.4}{12}\selectfont Richiami};
\node[anchor=north west,inner sep=0] at ([xshift=2.6cm,yshift=-10.64cm]current page.north west){\usebox{\RecallBox}};
\draw[Rule,line width=.35pt] ([xshift=.95cm,yshift=-11.92cm]current page.north west)--([xshift=19.05cm,yshift=-11.92cm]current page.north west);\node[anchor=north west,inner sep=0] at ([xshift=0.95cm,yshift=-12.06cm]current page.north west){\color{Muted}\fontsize{7.7}{9.3}\selectfont M13 · Cronologie e archivi sedimentari};
\node[anchor=north east,inner sep=0] at ([xshift=17.65cm,yshift=-12.06cm]current page.north west){\color{Classical}\fontsize{7.3}{9.1}\selectfont };
\node[anchor=north east,inner sep=0] at ([xshift=19.05cm,yshift=-12.06cm]current page.north west){\color{Muted}\fontsize{7.7}{9.3}\selectfont 3/30};
\end{tikzpicture}\end{frame}
\begin{lrbox}{\TitleBox}\begin{minipage}[t]{18.2cm}\raggedright\sloppy\color{Ink}\fontsize{20.7}{24}\selectfont \bfseries\color{Navy}L’associazione fra ancora ed evidenza stratigrafica\end{minipage}\end{lrbox}
\begin{lrbox}{\BodyBox}\begin{minipage}[t]{10.25cm}\raggedright\sloppy\color{Ink}\fontsize{12.4}{16.4}\selectfont Ogni ancora lega un’età a una profondità e a un insieme di metadati, che nelle indagini reali comprende materiale, metodo, identificativo e contesto deposizionale. Il controllo ordina le righe per profondità e verifica la coerenza delle età, mantenendo intatte le coppie. Ordinare separatamente le età altererebbe l’associazione con i livelli, cancellando l’evidenza di eventuali inversioni. La monotonicità può essere un’ipotesi geologica pertinente per una successione indisturbata, la cui validità deve essere valutata prima di imporla al modello.\par\vspace{.28cm}{\centering\fontsize{13}{17}\selectfont\(\displaystyle\begin{gathered}z_{i+1}>z_i,\qquad t_{i+1}>t_i\quad\text{nel modello monotono}\end{gathered}\)\par}\end{minipage}\end{lrbox}
\begin{lrbox}{\SideBox}\begin{minipage}[t]{7.45cm}\raggedright\sloppy\color{Ink}\fontsize{12.2}{16.1}\selectfont \input{tikz/M13/D01.tex}\end{minipage}\end{lrbox}
\begin{lrbox}{\RecallBox}\begin{minipage}[t]{16.45cm}\raggedright\sloppy\color{Ink}\fontsize{10.1}{13.2}\selectfont La coppia età–profondità costituisce l’unità informativa che il controllo deve preservare.\end{minipage}\end{lrbox}
\MeasureBox{4-TitleBox}{\TitleBox}
\MeasureBox{4-BodyBox}{\BodyBox}
\MeasureBox{4-SideBox}{\SideBox}
\MeasureBox{4-RecallBox}{\RecallBox}
\begin{frame}[plain]\hypertarget{slide4}{}
\bookmark[level=1,page=4]{L’associazione fra ancora ed evidenza stratigrafica}
\begin{tikzpicture}[remember picture,overlay]\node[anchor=north east,inner sep=0] at ([xshift=19.05cm,yshift=-0.16cm]current page.north west){\color{Classical}\fontsize{8}{10}\selectfont\hyperlink{slide2}{Indice}};
\node[anchor=north west,inner sep=0] at ([xshift=0.95cm,yshift=-0.62cm]current page.north west){\usebox{\TitleBox}};
\draw[Classical,line width=1.3pt] ([xshift=.95cm,yshift=-1.68cm]current page.north west)--([xshift=2.25cm,yshift=-1.68cm]current page.north west);\node[anchor=north west,inner sep=0] at ([xshift=0.95cm,yshift=-2.12cm]current page.north west){\usebox{\BodyBox}};
\node[anchor=north west,inner sep=0] at ([xshift=11.6cm,yshift=-2.1cm]current page.north west){\color{Quantum}\bfseries\fontsize{11.5}{14}\selectfont Lettura del diagramma};
\node[anchor=north west,inner sep=0] at ([xshift=11.6cm,yshift=-2.75cm]current page.north west){\usebox{\SideBox}};
\draw[Rule,line width=.5pt] ([xshift=.95cm,yshift=-10.43cm]current page.north west)--([xshift=19.05cm,yshift=-10.43cm]current page.north west);\node[anchor=north west,inner sep=0] at ([xshift=0.95cm,yshift=-10.66cm]current page.north west){\color{Classical}\bfseries\fontsize{9.4}{12}\selectfont Richiami};
\node[anchor=north west,inner sep=0] at ([xshift=2.6cm,yshift=-10.64cm]current page.north west){\usebox{\RecallBox}};
\draw[Rule,line width=.35pt] ([xshift=.95cm,yshift=-11.92cm]current page.north west)--([xshift=19.05cm,yshift=-11.92cm]current page.north west);\node[anchor=north west,inner sep=0] at ([xshift=0.95cm,yshift=-12.06cm]current page.north west){\color{Muted}\fontsize{7.7}{9.3}\selectfont M13 · Cronologie e archivi sedimentari};
\node[anchor=north east,inner sep=0] at ([xshift=17.65cm,yshift=-12.06cm]current page.north west){\color{Classical}\fontsize{7.3}{9.1}\selectfont };
\node[anchor=north east,inner sep=0] at ([xshift=19.05cm,yshift=-12.06cm]current page.north west){\color{Muted}\fontsize{7.7}{9.3}\selectfont 4/30};
\end{tikzpicture}\end{frame}
\begin{lrbox}{\TitleBox}\begin{minipage}[t]{18.2cm}\raggedright\sloppy\color{Ink}\fontsize{20.7}{24}\selectfont \bfseries\color{Navy}Origine BP e rappresentazione dell’anno civile\end{minipage}\end{lrbox}
\begin{lrbox}{\BodyBox}\begin{minipage}[t]{10.25cm}\raggedright\sloppy\color{Ink}\fontsize{12.4}{16.4}\selectfont La convenzione BP fissa il presente al 1950, indipendentemente dall’anno nel quale si esegue l’analisi. Per età già espresse in anni di calendario BP, la sottrazione da 1950 produce una coordinata sulla scala degli anni astronomici. Il passaggio alle etichette a.C. e d.C. richiede attenzione all’assenza dell’anno zero nella notazione storica: l’anno astronomico zero corrisponde al primo anno avanti Cristo. Le età negative in cal BP individuerebbero invece tempi successivi al riferimento convenzionale.\par\vspace{.28cm}{\centering\fontsize{13}{17}\selectfont\(\displaystyle\begin{gathered}y_{\rm astronomico}=1950-t_{\rm cal\ BP}\end{gathered}\)\par}\end{minipage}\end{lrbox}
\begin{lrbox}{\SideBox}\begin{minipage}[t]{7.45cm}\raggedright\sloppy\color{Ink}\fontsize{12.2}{16.1}\selectfont \input{tikz/M13/D02.tex}\end{minipage}\end{lrbox}
\begin{lrbox}{\RecallBox}\begin{minipage}[t]{16.45cm}\raggedright\sloppy\color{Ink}\fontsize{10.1}{13.2}\selectfont L’origine convenzionale e la notazione dell’anno devono accompagnare ogni cronologia.\end{minipage}\end{lrbox}
\MeasureBox{5-TitleBox}{\TitleBox}
\MeasureBox{5-BodyBox}{\BodyBox}
\MeasureBox{5-SideBox}{\SideBox}
\MeasureBox{5-RecallBox}{\RecallBox}
\begin{frame}[plain]\hypertarget{slide5}{}
\bookmark[level=1,page=5]{Origine BP e rappresentazione dell’anno civile}
\begin{tikzpicture}[remember picture,overlay]\node[anchor=north east,inner sep=0] at ([xshift=19.05cm,yshift=-0.16cm]current page.north west){\color{Classical}\fontsize{8}{10}\selectfont\hyperlink{slide2}{Indice}};
\node[anchor=north west,inner sep=0] at ([xshift=0.95cm,yshift=-0.62cm]current page.north west){\usebox{\TitleBox}};
\draw[Classical,line width=1.3pt] ([xshift=.95cm,yshift=-1.68cm]current page.north west)--([xshift=2.25cm,yshift=-1.68cm]current page.north west);\node[anchor=north west,inner sep=0] at ([xshift=0.95cm,yshift=-2.12cm]current page.north west){\usebox{\BodyBox}};
\node[anchor=north west,inner sep=0] at ([xshift=11.6cm,yshift=-2.1cm]current page.north west){\color{Quantum}\bfseries\fontsize{11.5}{14}\selectfont Lettura del diagramma};
\node[anchor=north west,inner sep=0] at ([xshift=11.6cm,yshift=-2.75cm]current page.north west){\usebox{\SideBox}};
\draw[Rule,line width=.5pt] ([xshift=.95cm,yshift=-10.43cm]current page.north west)--([xshift=19.05cm,yshift=-10.43cm]current page.north west);\node[anchor=north west,inner sep=0] at ([xshift=0.95cm,yshift=-10.66cm]current page.north west){\color{Classical}\bfseries\fontsize{9.4}{12}\selectfont Richiami};
\node[anchor=north west,inner sep=0] at ([xshift=2.6cm,yshift=-10.64cm]current page.north west){\usebox{\RecallBox}};
\draw[Rule,line width=.35pt] ([xshift=.95cm,yshift=-11.92cm]current page.north west)--([xshift=19.05cm,yshift=-11.92cm]current page.north west);\node[anchor=north west,inner sep=0] at ([xshift=0.95cm,yshift=-12.06cm]current page.north west){\color{Muted}\fontsize{7.7}{9.3}\selectfont M13 · Cronologie e archivi sedimentari};
\node[anchor=north east,inner sep=0] at ([xshift=17.65cm,yshift=-12.06cm]current page.north west){\color{Classical}\fontsize{7.3}{9.1}\selectfont IntCal20: convenzioni temporali};
\node[anchor=north east,inner sep=0] at ([xshift=19.05cm,yshift=-12.06cm]current page.north west){\color{Muted}\fontsize{7.7}{9.3}\selectfont 5/30};
\end{tikzpicture}\end{frame}
\begin{lrbox}{\TitleBox}\begin{minipage}[t]{18.2cm}\raggedright\sloppy\color{Ink}\fontsize{20.7}{24}\selectfont \bfseries\color{Navy}Età radiocarbonica e calibrazione\end{minipage}\end{lrbox}
\begin{lrbox}{\BodyBox}\begin{minipage}[t]{10.25cm}\raggedright\sloppy\color{Ink}\fontsize{12.4}{16.4}\selectfont Un’età radiocarbonica convenzionale richiede una calibrazione per essere espressa sulla scala del calendario, poiché la relazione fra le due scale dipende dalla storia del radiocarbonio e dal contesto del campione. Una semplice sottrazione da 1950 non esegue tale operazione. Curve di calibrazione, ambiente di formazione ed eventuali effetti di riserva devono essere scelti in modo pertinente. Le sei ancore del laboratorio sono già età di calendario inventate, per cui l’esercizio affronta la costruzione della cronologia senza simulare una calibrazione radiocarbonica reale.\end{minipage}\end{lrbox}
\begin{lrbox}{\SideBox}\begin{minipage}[t]{7.45cm}\raggedright\sloppy\color{Ink}\fontsize{12.2}{16.1}\selectfont Una distribuzione calibrata può essere asimmetrica o multimodale. Riassumerla automaticamente mediante una media e una deviazione standard può perdere informazione utile alla costruzione del modello stratigrafico. I dati reali richiedono la conservazione del risultato di calibrazione e delle sue assunzioni.\end{minipage}\end{lrbox}
\begin{lrbox}{\RecallBox}\begin{minipage}[t]{16.45cm}\raggedright\sloppy\color{Ink}\fontsize{10.1}{13.2}\selectfont Calibrazione della datazione e interpolazione età–profondità sono operazioni scientifiche distinte.\end{minipage}\end{lrbox}
\MeasureBox{6-TitleBox}{\TitleBox}
\MeasureBox{6-BodyBox}{\BodyBox}
\MeasureBox{6-SideBox}{\SideBox}
\MeasureBox{6-RecallBox}{\RecallBox}
\begin{frame}[plain]\hypertarget{slide6}{}
\bookmark[level=1,page=6]{Età radiocarbonica e calibrazione}
\begin{tikzpicture}[remember picture,overlay]\node[anchor=north east,inner sep=0] at ([xshift=19.05cm,yshift=-0.16cm]current page.north west){\color{Classical}\fontsize{8}{10}\selectfont\hyperlink{slide2}{Indice}};
\node[anchor=north west,inner sep=0] at ([xshift=0.95cm,yshift=-0.62cm]current page.north west){\usebox{\TitleBox}};
\draw[Classical,line width=1.3pt] ([xshift=.95cm,yshift=-1.68cm]current page.north west)--([xshift=2.25cm,yshift=-1.68cm]current page.north west);\node[anchor=north west,inner sep=0] at ([xshift=0.95cm,yshift=-2.12cm]current page.north west){\usebox{\BodyBox}};
\node[anchor=north west,inner sep=0] at ([xshift=11.6cm,yshift=-2.1cm]current page.north west){\color{Quantum}\bfseries\fontsize{11.5}{14}\selectfont Sviluppo e interpretazione};
\node[anchor=north west,inner sep=0] at ([xshift=11.6cm,yshift=-2.75cm]current page.north west){\usebox{\SideBox}};
\draw[Rule,line width=.5pt] ([xshift=.95cm,yshift=-10.43cm]current page.north west)--([xshift=19.05cm,yshift=-10.43cm]current page.north west);\node[anchor=north west,inner sep=0] at ([xshift=0.95cm,yshift=-10.66cm]current page.north west){\color{Classical}\bfseries\fontsize{9.4}{12}\selectfont Richiami};
\node[anchor=north west,inner sep=0] at ([xshift=2.6cm,yshift=-10.64cm]current page.north west){\usebox{\RecallBox}};
\draw[Rule,line width=.35pt] ([xshift=.95cm,yshift=-11.92cm]current page.north west)--([xshift=19.05cm,yshift=-11.92cm]current page.north west);\node[anchor=north west,inner sep=0] at ([xshift=0.95cm,yshift=-12.06cm]current page.north west){\color{Muted}\fontsize{7.7}{9.3}\selectfont M13 · Cronologie e archivi sedimentari};
\node[anchor=north east,inner sep=0] at ([xshift=17.65cm,yshift=-12.06cm]current page.north west){\color{Classical}\fontsize{7.3}{9.1}\selectfont Reimer et al., 2020; Heaton et al., 2020};
\node[anchor=north east,inner sep=0] at ([xshift=19.05cm,yshift=-12.06cm]current page.north west){\color{Muted}\fontsize{7.7}{9.3}\selectfont 6/30};
\end{tikzpicture}\end{frame}
\begin{lrbox}{\TitleBox}\begin{minipage}[t]{18.2cm}\raggedright\sloppy\color{Ink}\fontsize{20.7}{24}\selectfont \bfseries\color{Navy}Interpolazione lineare fra due ancore\end{minipage}\end{lrbox}
\begin{lrbox}{\BodyBox}\begin{minipage}[t]{10.25cm}\raggedright\sloppy\color{Ink}\fontsize{12.4}{16.4}\selectfont Fra due profondità adiacenti, l’interpolazione lineare combina le rispettive età con pesi determinati dalla posizione relativa. Il modello assume che il tasso apparente resti costante entro quel segmento e consente che cambi attraversando un’ancora. A 35 centimetri, posizione intermedia fra 20 e 50 centimetri, i pesi sono entrambi un mezzo e l’età risulta pari a 305 anni cal BP. Il passaggio esatto per le ancore deriva dalla definizione dell’interpolatore e non costituisce una verifica indipendente della cronologia.\par\vspace{.28cm}{\centering\fontsize{13}{17}\selectfont\(\displaystyle\begin{gathered}\lambda=\frac{z-z_i}{z_{i+1}-z_i},\qquad t(z)=(1-\lambda)t_i+\lambda t_{i+1}\end{gathered}\)\par}\end{minipage}\end{lrbox}
\begin{lrbox}{\SideBox}\begin{minipage}[t]{7.45cm}\raggedright\sloppy\color{Ink}\fontsize{12.2}{16.1}\selectfont \input{tikz/M13/D03.tex}\par\vspace{.16cm}{\fontsize{9.5}{12.6}\selectfont $z_i,t_i$: profondità ed età dell’ancora i; $\lambda$: peso dell’ancora successiva. Esempio in cm e anni cal BP, riferiti al 1950.}\end{minipage}\end{lrbox}
\begin{lrbox}{\RecallBox}\begin{minipage}[t]{16.45cm}\raggedright\sloppy\color{Ink}\fontsize{10.1}{13.2}\selectfont L’interpolazione incorpora un’ipotesi sul comportamento fra le ancore.\end{minipage}\end{lrbox}
\MeasureBox{7-TitleBox}{\TitleBox}
\MeasureBox{7-BodyBox}{\BodyBox}
\MeasureBox{7-SideBox}{\SideBox}
\MeasureBox{7-RecallBox}{\RecallBox}
\begin{frame}[plain]\hypertarget{slide7}{}
\bookmark[level=0,page=7]{Interpolazione e tassi di accumulo}
\bookmark[level=1,page=7]{Interpolazione lineare fra due ancore}
\begin{tikzpicture}[remember picture,overlay]\node[anchor=north east,inner sep=0] at ([xshift=19.05cm,yshift=-0.16cm]current page.north west){\color{Classical}\fontsize{8}{10}\selectfont\hyperlink{slide2}{Indice}};
\node[anchor=north west,inner sep=0] at ([xshift=0.95cm,yshift=-0.62cm]current page.north west){\usebox{\TitleBox}};
\draw[Classical,line width=1.3pt] ([xshift=.95cm,yshift=-1.68cm]current page.north west)--([xshift=2.25cm,yshift=-1.68cm]current page.north west);\node[anchor=north west,inner sep=0] at ([xshift=0.95cm,yshift=-2.12cm]current page.north west){\usebox{\BodyBox}};
\node[anchor=north west,inner sep=0] at ([xshift=11.6cm,yshift=-2.1cm]current page.north west){\color{Quantum}\bfseries\fontsize{11.5}{14}\selectfont Lettura del diagramma};
\node[anchor=north west,inner sep=0] at ([xshift=11.6cm,yshift=-2.75cm]current page.north west){\usebox{\SideBox}};
\draw[Rule,line width=.5pt] ([xshift=.95cm,yshift=-10.43cm]current page.north west)--([xshift=19.05cm,yshift=-10.43cm]current page.north west);\node[anchor=north west,inner sep=0] at ([xshift=0.95cm,yshift=-10.66cm]current page.north west){\color{Classical}\bfseries\fontsize{9.4}{12}\selectfont Richiami};
\node[anchor=north west,inner sep=0] at ([xshift=2.6cm,yshift=-10.64cm]current page.north west){\usebox{\RecallBox}};
\draw[Rule,line width=.35pt] ([xshift=.95cm,yshift=-11.92cm]current page.north west)--([xshift=19.05cm,yshift=-11.92cm]current page.north west);\node[anchor=north west,inner sep=0] at ([xshift=0.95cm,yshift=-12.06cm]current page.north west){\color{Muted}\fontsize{7.7}{9.3}\selectfont M13 · Cronologie e archivi sedimentari};
\node[anchor=north east,inner sep=0] at ([xshift=17.65cm,yshift=-12.06cm]current page.north west){\color{Classical}\fontsize{7.3}{9.1}\selectfont };
\node[anchor=north east,inner sep=0] at ([xshift=19.05cm,yshift=-12.06cm]current page.north west){\color{Muted}\fontsize{7.7}{9.3}\selectfont 7/30};
\end{tikzpicture}\end{frame}
\begin{lrbox}{\TitleBox}\begin{minipage}[t]{18.2cm}\raggedright\sloppy\color{Ink}\fontsize{20.7}{24}\selectfont \bfseries\color{Navy}Il dominio dell’interpolazione in R\end{minipage}\end{lrbox}
\begin{lrbox}{\BodyBox}\begin{minipage}[t]{10.25cm}\raggedright\sloppy\color{Ink}\fontsize{12.4}{16.4}\selectfont La funzione approx restituisce valori entro il dominio delimitato dalle ancore, mentre rule uguale a uno rende esplicita la mancanza di una previsione esterna. Stimare l’età a 180 centimetri richiederebbe un’ipotesi aggiuntiva sul tasso oltre l’ultima ancora. Prolungare l’ultimo segmento oppure mantenere costante l’età estrema definisce due modelli differenti, che devono essere motivati separatamente. Il controllo manuale a 35 centimetri e la verifica dei valori esterni costituiscono prove semplici della coerenza dell’implementazione.\end{minipage}\end{lrbox}
\begin{lrbox}{\SideBox}\begin{minipage}[t]{7.45cm}\raggedright\sloppy\color{Ink}\fontsize{12.2}{16.1}\selectfont \par\vspace{.25cm}\lstinputlisting[language=R]{estratti/M13/interpolazione.R}\end{minipage}\end{lrbox}
\begin{lrbox}{\RecallBox}\begin{minipage}[t]{16.45cm}\raggedright\sloppy\color{Ink}\fontsize{10.1}{13.2}\selectfont Ogni estrapolazione introduce ipotesi ulteriori rispetto all’interpolazione nel dominio osservato.\end{minipage}\end{lrbox}
\MeasureBox{8-TitleBox}{\TitleBox}
\MeasureBox{8-BodyBox}{\BodyBox}
\MeasureBox{8-SideBox}{\SideBox}
\MeasureBox{8-RecallBox}{\RecallBox}
\begin{frame}[plain]\hypertarget{slide8}{}
\bookmark[level=1,page=8]{Il dominio dell’interpolazione in R}
\begin{tikzpicture}[remember picture,overlay]\node[anchor=north east,inner sep=0] at ([xshift=19.05cm,yshift=-0.16cm]current page.north west){\color{Classical}\fontsize{8}{10}\selectfont\hyperlink{slide2}{Indice}};
\node[anchor=north west,inner sep=0] at ([xshift=0.95cm,yshift=-0.62cm]current page.north west){\usebox{\TitleBox}};
\draw[Classical,line width=1.3pt] ([xshift=.95cm,yshift=-1.68cm]current page.north west)--([xshift=2.25cm,yshift=-1.68cm]current page.north west);\node[anchor=north west,inner sep=0] at ([xshift=0.95cm,yshift=-2.12cm]current page.north west){\usebox{\BodyBox}};
\node[anchor=north west,inner sep=0] at ([xshift=11.6cm,yshift=-2.1cm]current page.north west){\color{Quantum}\bfseries\fontsize{11.5}{14}\selectfont Implementazione in R};
\node[anchor=north west,inner sep=0] at ([xshift=11.6cm,yshift=-2.75cm]current page.north west){\usebox{\SideBox}};
\draw[Rule,line width=.5pt] ([xshift=.95cm,yshift=-10.43cm]current page.north west)--([xshift=19.05cm,yshift=-10.43cm]current page.north west);\node[anchor=north west,inner sep=0] at ([xshift=0.95cm,yshift=-10.66cm]current page.north west){\color{Classical}\bfseries\fontsize{9.4}{12}\selectfont Richiami};
\node[anchor=north west,inner sep=0] at ([xshift=2.6cm,yshift=-10.64cm]current page.north west){\usebox{\RecallBox}};
\draw[Rule,line width=.35pt] ([xshift=.95cm,yshift=-11.92cm]current page.north west)--([xshift=19.05cm,yshift=-11.92cm]current page.north west);\node[anchor=north west,inner sep=0] at ([xshift=0.95cm,yshift=-12.06cm]current page.north west){\color{Muted}\fontsize{7.7}{9.3}\selectfont M13 · Cronologie e archivi sedimentari};
\node[anchor=north east,inner sep=0] at ([xshift=17.65cm,yshift=-12.06cm]current page.north west){\color{Classical}\fontsize{7.3}{9.1}\selectfont R: approx};
\node[anchor=north east,inner sep=0] at ([xshift=19.05cm,yshift=-12.06cm]current page.north west){\color{Muted}\fontsize{7.7}{9.3}\selectfont 8/30};
\end{tikzpicture}\end{frame}
\begin{lrbox}{\TitleBox}\begin{minipage}[t]{18.2cm}\raggedright\sloppy\color{Ink}\fontsize{20.7}{24}\selectfont \bfseries\color{Navy}Tasso apparente di accumulo e rapporto inverso\end{minipage}\end{lrbox}
\begin{lrbox}{\BodyBox}\begin{minipage}[t]{10.25cm}\raggedright\sloppy\color{Ink}\fontsize{12.4}{16.4}\selectfont Il rapporto fra incremento di profondità e intervallo temporale esprime un tasso apparente di accumulo, mentre il suo inverso misura anni per centimetro. Il termine apparente ricorda che spessore attuale, compattazione, erosione e lacune possono rendere il rapporto diverso da un flusso deposizionale originario. Nella successione sintetica tutti i segmenti hanno incrementi positivi. La conversione da centimetri per anno a millimetri per anno moltiplica per dieci e deve essere applicata ai valori prima di modificarne l’etichetta.\par\vspace{.28cm}{\centering\fontsize{13}{17}\selectfont\(\displaystyle\begin{gathered}r_i=\frac{z_{i+1}-z_i}{t_{i+1}-t_i},\qquad\frac{\mathrm dt}{\mathrm dz}=\frac1{r_i}\end{gathered}\)\par}\end{minipage}\end{lrbox}
\begin{lrbox}{\SideBox}\begin{minipage}[t]{7.45cm}\raggedright\sloppy\color{Ink}\fontsize{12.2}{16.1}\selectfont \includegraphics[width=7.45cm]{figure/M13/F02.pdf}\par\vspace{.16cm}{\fontsize{9.5}{12.6}\selectfont z in cm, t in anni: $r_i$ è in cm/anno. L’asse mostra $10r_i$ in mm/anno; le barre seguono gli intervalli di profondità.}\end{minipage}\end{lrbox}
\begin{lrbox}{\RecallBox}\begin{minipage}[t]{16.45cm}\raggedright\sloppy\color{Ink}\fontsize{10.1}{13.2}\selectfont Il tasso deve essere accompagnato dall’unità e dalle condizioni stratigrafiche della sua interpretazione.\end{minipage}\end{lrbox}
\MeasureBox{9-TitleBox}{\TitleBox}
\MeasureBox{9-BodyBox}{\BodyBox}
\MeasureBox{9-SideBox}{\SideBox}
\MeasureBox{9-RecallBox}{\RecallBox}
\begin{frame}[plain]\hypertarget{slide9}{}
\bookmark[level=1,page=9]{Tasso apparente di accumulo e rapporto inverso}
\begin{tikzpicture}[remember picture,overlay]\node[anchor=north east,inner sep=0] at ([xshift=19.05cm,yshift=-0.16cm]current page.north west){\color{Classical}\fontsize{8}{10}\selectfont\hyperlink{slide2}{Indice}};
\node[anchor=north west,inner sep=0] at ([xshift=0.95cm,yshift=-0.62cm]current page.north west){\usebox{\TitleBox}};
\draw[Classical,line width=1.3pt] ([xshift=.95cm,yshift=-1.68cm]current page.north west)--([xshift=2.25cm,yshift=-1.68cm]current page.north west);\node[anchor=north west,inner sep=0] at ([xshift=0.95cm,yshift=-2.12cm]current page.north west){\usebox{\BodyBox}};
\node[anchor=north west,inner sep=0] at ([xshift=11.6cm,yshift=-2.1cm]current page.north west){\color{Quantum}\bfseries\fontsize{11.5}{14}\selectfont Evidenza grafica};
\node[anchor=north west,inner sep=0] at ([xshift=11.6cm,yshift=-2.75cm]current page.north west){\usebox{\SideBox}};
\draw[Rule,line width=.5pt] ([xshift=.95cm,yshift=-10.43cm]current page.north west)--([xshift=19.05cm,yshift=-10.43cm]current page.north west);\node[anchor=north west,inner sep=0] at ([xshift=0.95cm,yshift=-10.66cm]current page.north west){\color{Classical}\bfseries\fontsize{9.4}{12}\selectfont Richiami};
\node[anchor=north west,inner sep=0] at ([xshift=2.6cm,yshift=-10.64cm]current page.north west){\usebox{\RecallBox}};
\draw[Rule,line width=.35pt] ([xshift=.95cm,yshift=-11.92cm]current page.north west)--([xshift=19.05cm,yshift=-11.92cm]current page.north west);\node[anchor=north west,inner sep=0] at ([xshift=0.95cm,yshift=-12.06cm]current page.north west){\color{Muted}\fontsize{7.7}{9.3}\selectfont M13 · Cronologie e archivi sedimentari};
\node[anchor=north east,inner sep=0] at ([xshift=17.65cm,yshift=-12.06cm]current page.north west){\color{Classical}\fontsize{7.3}{9.1}\selectfont };
\node[anchor=north east,inner sep=0] at ([xshift=19.05cm,yshift=-12.06cm]current page.north west){\color{Muted}\fontsize{7.7}{9.3}\selectfont 9/30};
\end{tikzpicture}\end{frame}
\begin{lrbox}{\TitleBox}\begin{minipage}[t]{18.2cm}\raggedright\sloppy\color{Ink}\fontsize{20.7}{24}\selectfont \bfseries\color{Navy}Media dei tassi e ponderazione temporale\end{minipage}\end{lrbox}
\begin{lrbox}{\BodyBox}\begin{minipage}[t]{10.25cm}\raggedright\sloppy\color{Ink}\fontsize{12.4}{16.4}\selectfont La media semplice dei cinque tassi assegna uguale peso a ogni segmento, benché i segmenti abbiano durate differenti. Per descrivere l’accumulo medio lungo l’intero periodo si ponderano invece i tassi per le rispettive durate; il risultato coincide con il rapporto fra spessore totale e durata complessiva. Nel caso didattico, 160 centimetri in 2200 anni corrispondono a circa 0,727 millimetri per anno. Il confronto rende esplicita la popolazione implicita nella media: segmenti equiponderati oppure tempo trascorso.\par\vspace{.28cm}{\centering\fontsize{13}{17}\selectfont\(\displaystyle\begin{gathered}\overline r_{\Delta t}=\frac{\sum_i\Delta t_i r_i}{\sum_i\Delta t_i}=\frac{\sum_i\Delta z_i}{\sum_i\Delta t_i}\end{gathered}\)\par}\end{minipage}\end{lrbox}
\begin{lrbox}{\SideBox}\begin{minipage}[t]{7.45cm}\raggedright\sloppy\color{Ink}\fontsize{12.2}{16.1}\selectfont \input{tikz/M13/D04.tex}\end{minipage}\end{lrbox}
\begin{lrbox}{\RecallBox}\begin{minipage}[t]{16.45cm}\raggedright\sloppy\color{Ink}\fontsize{10.1}{13.2}\selectfont La ponderazione deriva dalla domanda alla quale deve rispondere il tasso medio.\end{minipage}\end{lrbox}
\MeasureBox{10-TitleBox}{\TitleBox}
\MeasureBox{10-BodyBox}{\BodyBox}
\MeasureBox{10-SideBox}{\SideBox}
\MeasureBox{10-RecallBox}{\RecallBox}
\begin{frame}[plain]\hypertarget{slide10}{}
\bookmark[level=1,page=10]{Media dei tassi e ponderazione temporale}
\begin{tikzpicture}[remember picture,overlay]\node[anchor=north east,inner sep=0] at ([xshift=19.05cm,yshift=-0.16cm]current page.north west){\color{Classical}\fontsize{8}{10}\selectfont\hyperlink{slide2}{Indice}};
\node[anchor=north west,inner sep=0] at ([xshift=0.95cm,yshift=-0.62cm]current page.north west){\usebox{\TitleBox}};
\draw[Classical,line width=1.3pt] ([xshift=.95cm,yshift=-1.68cm]current page.north west)--([xshift=2.25cm,yshift=-1.68cm]current page.north west);\node[anchor=north west,inner sep=0] at ([xshift=0.95cm,yshift=-2.12cm]current page.north west){\usebox{\BodyBox}};
\node[anchor=north west,inner sep=0] at ([xshift=11.6cm,yshift=-2.1cm]current page.north west){\color{Quantum}\bfseries\fontsize{11.5}{14}\selectfont Lettura del diagramma};
\node[anchor=north west,inner sep=0] at ([xshift=11.6cm,yshift=-2.75cm]current page.north west){\usebox{\SideBox}};
\draw[Rule,line width=.5pt] ([xshift=.95cm,yshift=-10.43cm]current page.north west)--([xshift=19.05cm,yshift=-10.43cm]current page.north west);\node[anchor=north west,inner sep=0] at ([xshift=0.95cm,yshift=-10.66cm]current page.north west){\color{Classical}\bfseries\fontsize{9.4}{12}\selectfont Richiami};
\node[anchor=north west,inner sep=0] at ([xshift=2.6cm,yshift=-10.64cm]current page.north west){\usebox{\RecallBox}};
\draw[Rule,line width=.35pt] ([xshift=.95cm,yshift=-11.92cm]current page.north west)--([xshift=19.05cm,yshift=-11.92cm]current page.north west);\node[anchor=north west,inner sep=0] at ([xshift=0.95cm,yshift=-12.06cm]current page.north west){\color{Muted}\fontsize{7.7}{9.3}\selectfont M13 · Cronologie e archivi sedimentari};
\node[anchor=north east,inner sep=0] at ([xshift=17.65cm,yshift=-12.06cm]current page.north west){\color{Classical}\fontsize{7.3}{9.1}\selectfont };
\node[anchor=north east,inner sep=0] at ([xshift=19.05cm,yshift=-12.06cm]current page.north west){\color{Muted}\fontsize{7.7}{9.3}\selectfont 10/30};
\end{tikzpicture}\end{frame}
\begin{lrbox}{\TitleBox}\begin{minipage}[t]{18.2cm}\raggedright\sloppy\color{Ink}\fontsize{20.7}{24}\selectfont \bfseries\color{Navy}Perturbazioni delle ancore e modello di sensibilità\end{minipage}\end{lrbox}
\begin{lrbox}{\BodyBox}\begin{minipage}[t]{10.25cm}\raggedright\sloppy\color{Ink}\fontsize{12.4}{16.4}\selectfont Gli scenari aggiungono alle età di base perturbazioni gaussiane con deviazioni standard didattiche pari a 0, 25, 40, 55, 80 e 100 anni. Le perturbazioni delle diverse ancore sono assunte indipendenti prima dell’eventuale condizionamento alla monotonicità; la prima ancora rimane fissa. Questa costruzione quantifica la sensibilità alle ipotesi prescelte e non equivale alla stima di una cronologia bayesiana da datazioni reali. Seme e sequenza dei numeri casuali vengono conservati per rendere identico il confronto fra procedure.\par\vspace{.28cm}{\centering\fontsize{13}{17}\selectfont\(\displaystyle\begin{gathered}T_i^{(b)}=t_i+\sigma_i E_i^{(b)},\qquad E_i^{(b)}\overset{\rm iid}{\sim}N(0{,}1)\end{gathered}\)\par}\end{minipage}\end{lrbox}
\begin{lrbox}{\SideBox}\begin{minipage}[t]{7.45cm}\raggedright\sloppy\color{Ink}\fontsize{12.2}{16.1}\selectfont \input{tikz/M13/D05.tex}\par\vspace{.16cm}{\fontsize{9.5}{12.6}\selectfont i: ancora; b: scenario; $t_i$: età di base; $\sigma_i$: deviazione prescelta. Indipendenza prima del controllo di monotonicità.}\end{minipage}\end{lrbox}
\begin{lrbox}{\RecallBox}\begin{minipage}[t]{16.45cm}\raggedright\sloppy\color{Ink}\fontsize{10.1}{13.2}\selectfont La distribuzione degli scenari è definita dalle ipotesi sulle perturbazioni.\end{minipage}\end{lrbox}
\MeasureBox{11-TitleBox}{\TitleBox}
\MeasureBox{11-BodyBox}{\BodyBox}
\MeasureBox{11-SideBox}{\SideBox}
\MeasureBox{11-RecallBox}{\RecallBox}
\begin{frame}[plain]\hypertarget{slide11}{}
\bookmark[level=0,page=11]{Scenari cronologici e propagazione}
\bookmark[level=1,page=11]{Perturbazioni delle ancore e modello di sensibilità}
\begin{tikzpicture}[remember picture,overlay]\node[anchor=north east,inner sep=0] at ([xshift=19.05cm,yshift=-0.16cm]current page.north west){\color{Classical}\fontsize{8}{10}\selectfont\hyperlink{slide2}{Indice}};
\node[anchor=north west,inner sep=0] at ([xshift=0.95cm,yshift=-0.62cm]current page.north west){\usebox{\TitleBox}};
\draw[Classical,line width=1.3pt] ([xshift=.95cm,yshift=-1.68cm]current page.north west)--([xshift=2.25cm,yshift=-1.68cm]current page.north west);\node[anchor=north west,inner sep=0] at ([xshift=0.95cm,yshift=-2.12cm]current page.north west){\usebox{\BodyBox}};
\node[anchor=north west,inner sep=0] at ([xshift=11.6cm,yshift=-2.1cm]current page.north west){\color{Quantum}\bfseries\fontsize{11.5}{14}\selectfont Lettura del diagramma};
\node[anchor=north west,inner sep=0] at ([xshift=11.6cm,yshift=-2.75cm]current page.north west){\usebox{\SideBox}};
\draw[Rule,line width=.5pt] ([xshift=.95cm,yshift=-10.43cm]current page.north west)--([xshift=19.05cm,yshift=-10.43cm]current page.north west);\node[anchor=north west,inner sep=0] at ([xshift=0.95cm,yshift=-10.66cm]current page.north west){\color{Classical}\bfseries\fontsize{9.4}{12}\selectfont Richiami};
\node[anchor=north west,inner sep=0] at ([xshift=2.6cm,yshift=-10.64cm]current page.north west){\usebox{\RecallBox}};
\draw[Rule,line width=.35pt] ([xshift=.95cm,yshift=-11.92cm]current page.north west)--([xshift=19.05cm,yshift=-11.92cm]current page.north west);\node[anchor=north west,inner sep=0] at ([xshift=0.95cm,yshift=-12.06cm]current page.north west){\color{Muted}\fontsize{7.7}{9.3}\selectfont M13 · Cronologie e archivi sedimentari};
\node[anchor=north east,inner sep=0] at ([xshift=17.65cm,yshift=-12.06cm]current page.north west){\color{Classical}\fontsize{7.3}{9.1}\selectfont };
\node[anchor=north east,inner sep=0] at ([xshift=19.05cm,yshift=-12.06cm]current page.north west){\color{Muted}\fontsize{7.7}{9.3}\selectfont 11/30};
\end{tikzpicture}\end{frame}
\begin{lrbox}{\TitleBox}\begin{minipage}[t]{18.2cm}\raggedright\sloppy\color{Ink}\fontsize{20.7}{24}\selectfont \bfseries\color{Navy}Monotonicità e identità delle ancore\end{minipage}\end{lrbox}
\begin{lrbox}{\BodyBox}\begin{minipage}[t]{10.25cm}\raggedright\sloppy\color{Ink}\fontsize{12.4}{16.4}\selectfont Il controllo esamina ogni vettore di età mantenendo l’associazione con le profondità originarie. Uno scenario non monotono viene escluso esplicitamente, mentre ordinare le età lo trasformerebbe in una configurazione diversa da quella generata. Il rifiuto condiziona la distribuzione delle perturbazioni e può introdurre dipendenza fra ancore. Con i valori e la sequenza inclusi, tutti i cinquecento scenari superano il controllo; il codice mantiene comunque la verifica, che diventa sostanziale aumentando le deviazioni standard o riducendo la separazione temporale delle ancore.\end{minipage}\end{lrbox}
\begin{lrbox}{\SideBox}\begin{minipage}[t]{7.45cm}\raggedright\sloppy\color{Ink}\fontsize{12.2}{16.1}\selectfont \par\vspace{.25cm}\lstinputlisting[language=R]{estratti/M13/scenari.R}\end{minipage}\end{lrbox}
\begin{lrbox}{\RecallBox}\begin{minipage}[t]{16.45cm}\raggedright\sloppy\color{Ink}\fontsize{10.1}{13.2}\selectfont Un vincolo di ordine deve essere applicato senza trasferire un’età da una profondità all’altra.\end{minipage}\end{lrbox}
\MeasureBox{12-TitleBox}{\TitleBox}
\MeasureBox{12-BodyBox}{\BodyBox}
\MeasureBox{12-SideBox}{\SideBox}
\MeasureBox{12-RecallBox}{\RecallBox}
\begin{frame}[plain]\hypertarget{slide12}{}
\bookmark[level=1,page=12]{Monotonicità e identità delle ancore}
\begin{tikzpicture}[remember picture,overlay]\node[anchor=north east,inner sep=0] at ([xshift=19.05cm,yshift=-0.16cm]current page.north west){\color{Classical}\fontsize{8}{10}\selectfont\hyperlink{slide2}{Indice}};
\node[anchor=north west,inner sep=0] at ([xshift=0.95cm,yshift=-0.62cm]current page.north west){\usebox{\TitleBox}};
\draw[Classical,line width=1.3pt] ([xshift=.95cm,yshift=-1.68cm]current page.north west)--([xshift=2.25cm,yshift=-1.68cm]current page.north west);\node[anchor=north west,inner sep=0] at ([xshift=0.95cm,yshift=-2.12cm]current page.north west){\usebox{\BodyBox}};
\node[anchor=north west,inner sep=0] at ([xshift=11.6cm,yshift=-2.1cm]current page.north west){\color{Quantum}\bfseries\fontsize{11.5}{14}\selectfont Implementazione in R};
\node[anchor=north west,inner sep=0] at ([xshift=11.6cm,yshift=-2.75cm]current page.north west){\usebox{\SideBox}};
\draw[Rule,line width=.5pt] ([xshift=.95cm,yshift=-10.43cm]current page.north west)--([xshift=19.05cm,yshift=-10.43cm]current page.north west);\node[anchor=north west,inner sep=0] at ([xshift=0.95cm,yshift=-10.66cm]current page.north west){\color{Classical}\bfseries\fontsize{9.4}{12}\selectfont Richiami};
\node[anchor=north west,inner sep=0] at ([xshift=2.6cm,yshift=-10.64cm]current page.north west){\usebox{\RecallBox}};
\draw[Rule,line width=.35pt] ([xshift=.95cm,yshift=-11.92cm]current page.north west)--([xshift=19.05cm,yshift=-11.92cm]current page.north west);\node[anchor=north west,inner sep=0] at ([xshift=0.95cm,yshift=-12.06cm]current page.north west){\color{Muted}\fontsize{7.7}{9.3}\selectfont M13 · Cronologie e archivi sedimentari};
\node[anchor=north east,inner sep=0] at ([xshift=17.65cm,yshift=-12.06cm]current page.north west){\color{Classical}\fontsize{7.3}{9.1}\selectfont };
\node[anchor=north east,inner sep=0] at ([xshift=19.05cm,yshift=-12.06cm]current page.north west){\color{Muted}\fontsize{7.7}{9.3}\selectfont 12/30};
\end{tikzpicture}\end{frame}
\begin{lrbox}{\TitleBox}\begin{minipage}[t]{18.2cm}\raggedright\sloppy\color{Ink}\fontsize{20.7}{24}\selectfont \bfseries\color{Navy}Bande puntuali del modello età–profondità\end{minipage}\end{lrbox}
\begin{lrbox}{\BodyBox}\begin{minipage}[t]{10.25cm}\raggedright\sloppy\color{Ink}\fontsize{12.4}{16.4}\selectfont Ogni scenario accettato viene interpolato sulle medesime profondità e produce una curva cronologica completa. I quantili al 2,5\%, 50\% e 97,5\% vengono calcolati separatamente a ogni profondità, definendo una banda puntuale degli scenari. Tale banda esprime la variabilità indotta dalle perturbazioni specificate; non possiede automaticamente una copertura frequentista certificata né una probabilità simultanea per l’intera curva. L’incertezza sulla forma fra le ancore rimane esclusa quando tutti gli scenari utilizzano la medesima interpolazione lineare.\end{minipage}\end{lrbox}
\begin{lrbox}{\SideBox}\begin{minipage}[t]{7.45cm}\raggedright\sloppy\color{Ink}\fontsize{12.2}{16.1}\selectfont \includegraphics[width=7.45cm]{figure/M13/F03.pdf}\par\vspace{.16cm}{\fontsize{9.5}{12.6}\selectfont Mediana e quantili puntuali degli scenari monotoni. Le ancore di base sono indicate in viola.}\end{minipage}\end{lrbox}
\begin{lrbox}{\RecallBox}\begin{minipage}[t]{16.45cm}\raggedright\sloppy\color{Ink}\fontsize{10.1}{13.2}\selectfont Una banda deve dichiarare quali fonti di incertezza siano state propagate e quali restino escluse.\end{minipage}\end{lrbox}
\MeasureBox{13-TitleBox}{\TitleBox}
\MeasureBox{13-BodyBox}{\BodyBox}
\MeasureBox{13-SideBox}{\SideBox}
\MeasureBox{13-RecallBox}{\RecallBox}
\begin{frame}[plain]\hypertarget{slide13}{}
\bookmark[level=1,page=13]{Bande puntuali del modello età–profondità}
\begin{tikzpicture}[remember picture,overlay]\node[anchor=north east,inner sep=0] at ([xshift=19.05cm,yshift=-0.16cm]current page.north west){\color{Classical}\fontsize{8}{10}\selectfont\hyperlink{slide2}{Indice}};
\node[anchor=north west,inner sep=0] at ([xshift=0.95cm,yshift=-0.62cm]current page.north west){\usebox{\TitleBox}};
\draw[Classical,line width=1.3pt] ([xshift=.95cm,yshift=-1.68cm]current page.north west)--([xshift=2.25cm,yshift=-1.68cm]current page.north west);\node[anchor=north west,inner sep=0] at ([xshift=0.95cm,yshift=-2.12cm]current page.north west){\usebox{\BodyBox}};
\node[anchor=north west,inner sep=0] at ([xshift=11.6cm,yshift=-2.1cm]current page.north west){\color{Quantum}\bfseries\fontsize{11.5}{14}\selectfont Evidenza grafica};
\node[anchor=north west,inner sep=0] at ([xshift=11.6cm,yshift=-2.75cm]current page.north west){\usebox{\SideBox}};
\draw[Rule,line width=.5pt] ([xshift=.95cm,yshift=-10.43cm]current page.north west)--([xshift=19.05cm,yshift=-10.43cm]current page.north west);\node[anchor=north west,inner sep=0] at ([xshift=0.95cm,yshift=-10.66cm]current page.north west){\color{Classical}\bfseries\fontsize{9.4}{12}\selectfont Richiami};
\node[anchor=north west,inner sep=0] at ([xshift=2.6cm,yshift=-10.64cm]current page.north west){\usebox{\RecallBox}};
\draw[Rule,line width=.35pt] ([xshift=.95cm,yshift=-11.92cm]current page.north west)--([xshift=19.05cm,yshift=-11.92cm]current page.north west);\node[anchor=north west,inner sep=0] at ([xshift=0.95cm,yshift=-12.06cm]current page.north west){\color{Muted}\fontsize{7.7}{9.3}\selectfont M13 · Cronologie e archivi sedimentari};
\node[anchor=north east,inner sep=0] at ([xshift=17.65cm,yshift=-12.06cm]current page.north west){\color{Classical}\fontsize{7.3}{9.1}\selectfont };
\node[anchor=north east,inner sep=0] at ([xshift=19.05cm,yshift=-12.06cm]current page.north west){\color{Muted}\fontsize{7.7}{9.3}\selectfont 13/30};
\end{tikzpicture}\end{frame}
\begin{lrbox}{\TitleBox}\begin{minipage}[t]{18.2cm}\raggedright\sloppy\color{Ink}\fontsize{20.7}{24}\selectfont \bfseries\color{Navy}Distribuzione dell’età a una profondità prefissata\end{minipage}\end{lrbox}
\begin{lrbox}{\BodyBox}\begin{minipage}[t]{10.25cm}\raggedright\sloppy\color{Ink}\fontsize{12.4}{16.4}\selectfont A cento centimetri, la cronologia di base assegna 1090 anni cal BP, combinando con pesi uguali le ancore a ottanta e centoventi centimetri. La distribuzione degli scenari restituisce una mediana prossima a 1087,5 anni e quantili al 2,5\% e 97,5\% prossimi a 997,7 e 1186,5 anni. Le differenze fra sintesi Monte Carlo e valore di base riflettono la sequenza finita delle perturbazioni. Raddoppiare le deviazioni standard definirebbe una nuova analisi di sensibilità, della quale andrebbero riportati anche gli eventuali scenari esclusi.\par\vspace{.28cm}{\centering\fontsize{13}{17}\selectfont\(\displaystyle\begin{gathered}T(100)=\tfrac12T(80)+\tfrac12T(120)\end{gathered}\)\par}\end{minipage}\end{lrbox}
\begin{lrbox}{\SideBox}\begin{minipage}[t]{7.45cm}\raggedright\sloppy\color{Ink}\fontsize{12.2}{16.1}\selectfont \includegraphics[width=7.45cm]{figure/M13/F04.pdf}\par\vspace{.16cm}{\fontsize{9.5}{12.6}\selectfont Età a 100 centimetri nei cinquecento scenari. La linea viola indica la cronologia di base.}\end{minipage}\end{lrbox}
\begin{lrbox}{\RecallBox}\begin{minipage}[t]{16.45cm}\raggedright\sloppy\color{Ink}\fontsize{10.1}{13.2}\selectfont I quantili descrivono gli scenari adottati e non aggiungono nuove evidenze di datazione.\end{minipage}\end{lrbox}
\MeasureBox{14-TitleBox}{\TitleBox}
\MeasureBox{14-BodyBox}{\BodyBox}
\MeasureBox{14-SideBox}{\SideBox}
\MeasureBox{14-RecallBox}{\RecallBox}
\begin{frame}[plain]\hypertarget{slide14}{}
\bookmark[level=1,page=14]{Distribuzione dell’età a una profondità prefissata}
\begin{tikzpicture}[remember picture,overlay]\node[anchor=north east,inner sep=0] at ([xshift=19.05cm,yshift=-0.16cm]current page.north west){\color{Classical}\fontsize{8}{10}\selectfont\hyperlink{slide2}{Indice}};
\node[anchor=north west,inner sep=0] at ([xshift=0.95cm,yshift=-0.62cm]current page.north west){\usebox{\TitleBox}};
\draw[Classical,line width=1.3pt] ([xshift=.95cm,yshift=-1.68cm]current page.north west)--([xshift=2.25cm,yshift=-1.68cm]current page.north west);\node[anchor=north west,inner sep=0] at ([xshift=0.95cm,yshift=-2.12cm]current page.north west){\usebox{\BodyBox}};
\node[anchor=north west,inner sep=0] at ([xshift=11.6cm,yshift=-2.1cm]current page.north west){\color{Quantum}\bfseries\fontsize{11.5}{14}\selectfont Evidenza grafica};
\node[anchor=north west,inner sep=0] at ([xshift=11.6cm,yshift=-2.75cm]current page.north west){\usebox{\SideBox}};
\draw[Rule,line width=.5pt] ([xshift=.95cm,yshift=-10.43cm]current page.north west)--([xshift=19.05cm,yshift=-10.43cm]current page.north west);\node[anchor=north west,inner sep=0] at ([xshift=0.95cm,yshift=-10.66cm]current page.north west){\color{Classical}\bfseries\fontsize{9.4}{12}\selectfont Richiami};
\node[anchor=north west,inner sep=0] at ([xshift=2.6cm,yshift=-10.64cm]current page.north west){\usebox{\RecallBox}};
\draw[Rule,line width=.35pt] ([xshift=.95cm,yshift=-11.92cm]current page.north west)--([xshift=19.05cm,yshift=-11.92cm]current page.north west);\node[anchor=north west,inner sep=0] at ([xshift=0.95cm,yshift=-12.06cm]current page.north west){\color{Muted}\fontsize{7.7}{9.3}\selectfont M13 · Cronologie e archivi sedimentari};
\node[anchor=north east,inner sep=0] at ([xshift=17.65cm,yshift=-12.06cm]current page.north west){\color{Classical}\fontsize{7.3}{9.1}\selectfont };
\node[anchor=north east,inner sep=0] at ([xshift=19.05cm,yshift=-12.06cm]current page.north west){\color{Muted}\fontsize{7.7}{9.3}\selectfont 14/30};
\end{tikzpicture}\end{frame}
\begin{lrbox}{\TitleBox}\begin{minipage}[t]{18.2cm}\raggedright\sloppy\color{Ink}\fontsize{20.7}{24}\selectfont \bfseries\color{Navy}Propagazione lineare e covarianza fra età\end{minipage}\end{lrbox}
\begin{lrbox}{\BodyBox}\begin{minipage}[t]{10.25cm}\raggedright\sloppy\color{Ink}\fontsize{12.4}{16.4}\selectfont Entro la medesima specificazione lineare, l’età a una profondità può essere scritta come combinazione delle età delle ancore. La varianza si ricava mediante i pesi di interpolazione e la matrice di covarianza delle ancore. Due profondità possono condividere una o più ancore, producendo covarianza positiva anche quando le perturbazioni delle ancore sono inizialmente indipendenti. Il condizionamento alla monotonicità può modificare la covarianza, richiedendone un aggiornamento coerente.\par\vspace{.28cm}{\centering\fontsize{13}{17}\selectfont\(\displaystyle\begin{gathered}T(z)=w(z)^\top T,\qquad\operatorname{Var}[T(z)]=w(z)^\top\Sigma w(z)\\ \operatorname{Cov}[T(z_1),T(z_2)]=w(z_1)^\top\Sigma w(z_2)\end{gathered}\)\par}\end{minipage}\end{lrbox}
\begin{lrbox}{\SideBox}\begin{minipage}[t]{7.45cm}\raggedright\sloppy\color{Ink}\fontsize{12.2}{16.1}\selectfont \input{tikz/M13/D06.tex}\par\vspace{.16cm}{\fontsize{9.5}{12.6}\selectfont T: vettore delle sei età; w(z): pesi d’interpolazione alla profondità z; $\Sigma$: covarianza delle ancore. Le età derivate condividono i pesi.}\end{minipage}\end{lrbox}
\begin{lrbox}{\RecallBox}\begin{minipage}[t]{16.45cm}\raggedright\sloppy\color{Ink}\fontsize{10.1}{13.2}\selectfont La propagazione deve conservare la dipendenza indotta dalle ancore comuni.\end{minipage}\end{lrbox}
\MeasureBox{15-TitleBox}{\TitleBox}
\MeasureBox{15-BodyBox}{\BodyBox}
\MeasureBox{15-SideBox}{\SideBox}
\MeasureBox{15-RecallBox}{\RecallBox}
\begin{frame}[plain]\hypertarget{slide15}{}
\bookmark[level=1,page=15]{Propagazione lineare e covarianza fra età}
\begin{tikzpicture}[remember picture,overlay]\node[anchor=north east,inner sep=0] at ([xshift=19.05cm,yshift=-0.16cm]current page.north west){\color{Classical}\fontsize{8}{10}\selectfont\hyperlink{slide2}{Indice}};
\node[anchor=north west,inner sep=0] at ([xshift=0.95cm,yshift=-0.62cm]current page.north west){\usebox{\TitleBox}};
\draw[Classical,line width=1.3pt] ([xshift=.95cm,yshift=-1.68cm]current page.north west)--([xshift=2.25cm,yshift=-1.68cm]current page.north west);\node[anchor=north west,inner sep=0] at ([xshift=0.95cm,yshift=-2.12cm]current page.north west){\usebox{\BodyBox}};
\node[anchor=north west,inner sep=0] at ([xshift=11.6cm,yshift=-2.1cm]current page.north west){\color{Quantum}\bfseries\fontsize{11.5}{14}\selectfont Lettura del diagramma};
\node[anchor=north west,inner sep=0] at ([xshift=11.6cm,yshift=-2.75cm]current page.north west){\usebox{\SideBox}};
\draw[Rule,line width=.5pt] ([xshift=.95cm,yshift=-10.43cm]current page.north west)--([xshift=19.05cm,yshift=-10.43cm]current page.north west);\node[anchor=north west,inner sep=0] at ([xshift=0.95cm,yshift=-10.66cm]current page.north west){\color{Classical}\bfseries\fontsize{9.4}{12}\selectfont Richiami};
\node[anchor=north west,inner sep=0] at ([xshift=2.6cm,yshift=-10.64cm]current page.north west){\usebox{\RecallBox}};
\draw[Rule,line width=.35pt] ([xshift=.95cm,yshift=-11.92cm]current page.north west)--([xshift=19.05cm,yshift=-11.92cm]current page.north west);\node[anchor=north west,inner sep=0] at ([xshift=0.95cm,yshift=-12.06cm]current page.north west){\color{Muted}\fontsize{7.7}{9.3}\selectfont M13 · Cronologie e archivi sedimentari};
\node[anchor=north east,inner sep=0] at ([xshift=17.65cm,yshift=-12.06cm]current page.north west){\color{Classical}\fontsize{7.3}{9.1}\selectfont };
\node[anchor=north east,inner sep=0] at ([xshift=19.05cm,yshift=-12.06cm]current page.north west){\color{Muted}\fontsize{7.7}{9.3}\selectfont 15/30};
\end{tikzpicture}\end{frame}
\begin{lrbox}{\TitleBox}\begin{minipage}[t]{18.2cm}\raggedright\sloppy\color{Ink}\fontsize{20.7}{24}\selectfont \bfseries\color{Navy}Il proxy sulla scala della profondità\end{minipage}\end{lrbox}
\begin{lrbox}{\BodyBox}\begin{minipage}[t]{10.25cm}\raggedright\sloppy\color{Ink}\fontsize{12.4}{16.4}\selectfont Il proxy sintetico comprende ottantuno valori acquisiti idealmente a passo costante di due centimetri. Su questa scala, intervalli uguali dell’asse rappresentano spessori uguali, mentre oscillazioni e tendenze descrivono variazioni lungo la successione. La regolarità spaziale consente di riconoscere l’ordine dei campioni, ma lascia aperta la durata rappresentata da ciascun tratto. Prima di collegare un massimo o una variazione a un evento temporale occorre applicare la cronologia e valutarne la risoluzione, insieme al rumore del proxy e al supporto del campione.\end{minipage}\end{lrbox}
\begin{lrbox}{\SideBox}\begin{minipage}[t]{7.45cm}\raggedright\sloppy\color{Ink}\fontsize{12.2}{16.1}\selectfont \includegraphics[width=7.45cm]{figure/M13/F05.pdf}\par\vspace{.16cm}{\fontsize{9.5}{12.6}\selectfont Ottantuno osservazioni sintetiche, separate da due centimetri. I punti conservano la posizione dei campioni originari.}\end{minipage}\end{lrbox}
\begin{lrbox}{\RecallBox}\begin{minipage}[t]{16.45cm}\raggedright\sloppy\color{Ink}\fontsize{10.1}{13.2}\selectfont La scala della profondità descrive lo spessore dell’archivio e deve essere distinta dalla scala del tempo.\end{minipage}\end{lrbox}
\MeasureBox{16-TitleBox}{\TitleBox}
\MeasureBox{16-BodyBox}{\BodyBox}
\MeasureBox{16-SideBox}{\SideBox}
\MeasureBox{16-RecallBox}{\RecallBox}
\begin{frame}[plain]\hypertarget{slide16}{}
\bookmark[level=0,page=16]{Proxy e risoluzione temporale}
\bookmark[level=1,page=16]{Il proxy sulla scala della profondità}
\begin{tikzpicture}[remember picture,overlay]\node[anchor=north east,inner sep=0] at ([xshift=19.05cm,yshift=-0.16cm]current page.north west){\color{Classical}\fontsize{8}{10}\selectfont\hyperlink{slide2}{Indice}};
\node[anchor=north west,inner sep=0] at ([xshift=0.95cm,yshift=-0.62cm]current page.north west){\usebox{\TitleBox}};
\draw[Classical,line width=1.3pt] ([xshift=.95cm,yshift=-1.68cm]current page.north west)--([xshift=2.25cm,yshift=-1.68cm]current page.north west);\node[anchor=north west,inner sep=0] at ([xshift=0.95cm,yshift=-2.12cm]current page.north west){\usebox{\BodyBox}};
\node[anchor=north west,inner sep=0] at ([xshift=11.6cm,yshift=-2.1cm]current page.north west){\color{Quantum}\bfseries\fontsize{11.5}{14}\selectfont Evidenza grafica};
\node[anchor=north west,inner sep=0] at ([xshift=11.6cm,yshift=-2.75cm]current page.north west){\usebox{\SideBox}};
\draw[Rule,line width=.5pt] ([xshift=.95cm,yshift=-10.43cm]current page.north west)--([xshift=19.05cm,yshift=-10.43cm]current page.north west);\node[anchor=north west,inner sep=0] at ([xshift=0.95cm,yshift=-10.66cm]current page.north west){\color{Classical}\bfseries\fontsize{9.4}{12}\selectfont Richiami};
\node[anchor=north west,inner sep=0] at ([xshift=2.6cm,yshift=-10.64cm]current page.north west){\usebox{\RecallBox}};
\draw[Rule,line width=.35pt] ([xshift=.95cm,yshift=-11.92cm]current page.north west)--([xshift=19.05cm,yshift=-11.92cm]current page.north west);\node[anchor=north west,inner sep=0] at ([xshift=0.95cm,yshift=-12.06cm]current page.north west){\color{Muted}\fontsize{7.7}{9.3}\selectfont M13 · Cronologie e archivi sedimentari};
\node[anchor=north east,inner sep=0] at ([xshift=17.65cm,yshift=-12.06cm]current page.north west){\color{Classical}\fontsize{7.3}{9.1}\selectfont };
\node[anchor=north east,inner sep=0] at ([xshift=19.05cm,yshift=-12.06cm]current page.north west){\color{Muted}\fontsize{7.7}{9.3}\selectfont 16/30};
\end{tikzpicture}\end{frame}
\begin{lrbox}{\TitleBox}\begin{minipage}[t]{18.2cm}\raggedright\sloppy\color{Ink}\fontsize{20.7}{24}\selectfont \bfseries\color{Navy}Il medesimo proxy sulla scala del tempo\end{minipage}\end{lrbox}
\begin{lrbox}{\BodyBox}\begin{minipage}[t]{10.25cm}\raggedright\sloppy\color{Ink}\fontsize{12.4}{16.4}\selectfont La trasformazione dell’asse assegna un’età a ogni profondità, lasciando invariati i valori del proxy. Poiché i tassi apparenti cambiano fra i segmenti, intervalli uguali in profondità corrispondono a durate diverse e la densità temporale dei campioni varia. Il passo originario di 2 cm corrisponde a 16 anni nel primo segmento e a 40 nell’ultimo, dilatando diversamente i tratti dell’asse. Qualunque analisi che richieda campionamento regolare nel tempo deve quindi verificare il nuovo asse, anziché trasferire automaticamente la frequenza definita in centimetri.\end{minipage}\end{lrbox}
\begin{lrbox}{\SideBox}\begin{minipage}[t]{7.45cm}\raggedright\sloppy\color{Ink}\fontsize{12.2}{16.1}\selectfont \includegraphics[width=7.45cm]{figure/M13/F06.pdf}\par\vspace{.16cm}{\fontsize{9.5}{12.6}\selectfont Gli stessi ottantuno valori dopo la conversione mediante la cronologia di base. Le distanze temporali sono irregolari.}\end{minipage}\end{lrbox}
\begin{lrbox}{\RecallBox}\begin{minipage}[t]{16.45cm}\raggedright\sloppy\color{Ink}\fontsize{10.1}{13.2}\selectfont Trasformare l’asse modifica le distanze fra osservazioni e il significato delle procedure temporali.\end{minipage}\end{lrbox}
\MeasureBox{17-TitleBox}{\TitleBox}
\MeasureBox{17-BodyBox}{\BodyBox}
\MeasureBox{17-SideBox}{\SideBox}
\MeasureBox{17-RecallBox}{\RecallBox}
\begin{frame}[plain]\hypertarget{slide17}{}
\bookmark[level=1,page=17]{Il medesimo proxy sulla scala del tempo}
\begin{tikzpicture}[remember picture,overlay]\node[anchor=north east,inner sep=0] at ([xshift=19.05cm,yshift=-0.16cm]current page.north west){\color{Classical}\fontsize{8}{10}\selectfont\hyperlink{slide2}{Indice}};
\node[anchor=north west,inner sep=0] at ([xshift=0.95cm,yshift=-0.62cm]current page.north west){\usebox{\TitleBox}};
\draw[Classical,line width=1.3pt] ([xshift=.95cm,yshift=-1.68cm]current page.north west)--([xshift=2.25cm,yshift=-1.68cm]current page.north west);\node[anchor=north west,inner sep=0] at ([xshift=0.95cm,yshift=-2.12cm]current page.north west){\usebox{\BodyBox}};
\node[anchor=north west,inner sep=0] at ([xshift=11.6cm,yshift=-2.1cm]current page.north west){\color{Quantum}\bfseries\fontsize{11.5}{14}\selectfont Evidenza grafica};
\node[anchor=north west,inner sep=0] at ([xshift=11.6cm,yshift=-2.75cm]current page.north west){\usebox{\SideBox}};
\draw[Rule,line width=.5pt] ([xshift=.95cm,yshift=-10.43cm]current page.north west)--([xshift=19.05cm,yshift=-10.43cm]current page.north west);\node[anchor=north west,inner sep=0] at ([xshift=0.95cm,yshift=-10.66cm]current page.north west){\color{Classical}\bfseries\fontsize{9.4}{12}\selectfont Richiami};
\node[anchor=north west,inner sep=0] at ([xshift=2.6cm,yshift=-10.64cm]current page.north west){\usebox{\RecallBox}};
\draw[Rule,line width=.35pt] ([xshift=.95cm,yshift=-11.92cm]current page.north west)--([xshift=19.05cm,yshift=-11.92cm]current page.north west);\node[anchor=north west,inner sep=0] at ([xshift=0.95cm,yshift=-12.06cm]current page.north west){\color{Muted}\fontsize{7.7}{9.3}\selectfont M13 · Cronologie e archivi sedimentari};
\node[anchor=north east,inner sep=0] at ([xshift=17.65cm,yshift=-12.06cm]current page.north west){\color{Classical}\fontsize{7.3}{9.1}\selectfont };
\node[anchor=north east,inner sep=0] at ([xshift=19.05cm,yshift=-12.06cm]current page.north west){\color{Muted}\fontsize{7.7}{9.3}\selectfont 17/30};
\end{tikzpicture}\end{frame}
\begin{lrbox}{\TitleBox}\begin{minipage}[t]{18.2cm}\raggedright\sloppy\color{Ink}\fontsize{20.7}{24}\selectfont \bfseries\color{Navy}Passo di campionamento e durata effettiva\end{minipage}\end{lrbox}
\begin{lrbox}{\BodyBox}\begin{minipage}[t]{10.25cm}\raggedright\sloppy\color{Ink}\fontsize{12.4}{16.4}\selectfont La durata fra due campioni consecutivi dipende dalla pendenza locale della curva età–profondità. Nel laboratorio, il passo di due centimetri rappresenta sedici anni nel primo segmento e quaranta anni nell’ultimo. Una variazione breve può pertanto risultare diversamente risolta lungo la successione, anche se il protocollo di campionamento resta costante in profondità. L’interpolazione cronologica rende esplicita questa eterogeneità, alla quale si aggiungono incertezze delle ancore e durata integrata dal supporto fisico di ciascun campione.\par\vspace{.28cm}{\centering\fontsize{13}{17}\selectfont\(\displaystyle\begin{gathered}\Delta t_i=t(z_{i+1})-t(z_i)\end{gathered}\)\par}\end{minipage}\end{lrbox}
\begin{lrbox}{\SideBox}\begin{minipage}[t]{7.45cm}\raggedright\sloppy\color{Ink}\fontsize{12.2}{16.1}\selectfont \includegraphics[width=7.45cm]{figure/M13/F12.pdf}\par\vspace{.16cm}{\fontsize{9.5}{12.6}\selectfont Durata fra due campioni adiacenti, rappresentata alla profondità intermedia. Il passo spaziale resta pari a due centimetri.}\end{minipage}\end{lrbox}
\begin{lrbox}{\RecallBox}\begin{minipage}[t]{16.45cm}\raggedright\sloppy\color{Ink}\fontsize{10.1}{13.2}\selectfont La risoluzione temporale deve essere valutata localmente lungo l’archivio.\end{minipage}\end{lrbox}
\MeasureBox{18-TitleBox}{\TitleBox}
\MeasureBox{18-BodyBox}{\BodyBox}
\MeasureBox{18-SideBox}{\SideBox}
\MeasureBox{18-RecallBox}{\RecallBox}
\begin{frame}[plain]\hypertarget{slide18}{}
\bookmark[level=1,page=18]{Passo di campionamento e durata effettiva}
\begin{tikzpicture}[remember picture,overlay]\node[anchor=north east,inner sep=0] at ([xshift=19.05cm,yshift=-0.16cm]current page.north west){\color{Classical}\fontsize{8}{10}\selectfont\hyperlink{slide2}{Indice}};
\node[anchor=north west,inner sep=0] at ([xshift=0.95cm,yshift=-0.62cm]current page.north west){\usebox{\TitleBox}};
\draw[Classical,line width=1.3pt] ([xshift=.95cm,yshift=-1.68cm]current page.north west)--([xshift=2.25cm,yshift=-1.68cm]current page.north west);\node[anchor=north west,inner sep=0] at ([xshift=0.95cm,yshift=-2.12cm]current page.north west){\usebox{\BodyBox}};
\node[anchor=north west,inner sep=0] at ([xshift=11.6cm,yshift=-2.1cm]current page.north west){\color{Quantum}\bfseries\fontsize{11.5}{14}\selectfont Evidenza grafica};
\node[anchor=north west,inner sep=0] at ([xshift=11.6cm,yshift=-2.75cm]current page.north west){\usebox{\SideBox}};
\draw[Rule,line width=.5pt] ([xshift=.95cm,yshift=-10.43cm]current page.north west)--([xshift=19.05cm,yshift=-10.43cm]current page.north west);\node[anchor=north west,inner sep=0] at ([xshift=0.95cm,yshift=-10.66cm]current page.north west){\color{Classical}\bfseries\fontsize{9.4}{12}\selectfont Richiami};
\node[anchor=north west,inner sep=0] at ([xshift=2.6cm,yshift=-10.64cm]current page.north west){\usebox{\RecallBox}};
\draw[Rule,line width=.35pt] ([xshift=.95cm,yshift=-11.92cm]current page.north west)--([xshift=19.05cm,yshift=-11.92cm]current page.north west);\node[anchor=north west,inner sep=0] at ([xshift=0.95cm,yshift=-12.06cm]current page.north west){\color{Muted}\fontsize{7.7}{9.3}\selectfont M13 · Cronologie e archivi sedimentari};
\node[anchor=north east,inner sep=0] at ([xshift=17.65cm,yshift=-12.06cm]current page.north west){\color{Classical}\fontsize{7.3}{9.1}\selectfont };
\node[anchor=north east,inner sep=0] at ([xshift=19.05cm,yshift=-12.06cm]current page.north west){\color{Muted}\fontsize{7.7}{9.3}\selectfont 18/30};
\end{tikzpicture}\end{frame}
\begin{lrbox}{\TitleBox}\begin{minipage}[t]{18.2cm}\raggedright\sloppy\color{Ink}\fontsize{20.7}{24}\selectfont \bfseries\color{Navy}Finestre per indice e finestre temporali\end{minipage}\end{lrbox}
\begin{lrbox}{\BodyBox}\begin{minipage}[t]{10.25cm}\raggedright\sloppy\color{Ink}\fontsize{12.4}{16.4}\selectfont Una finestra di cinque campioni copre sempre otto centimetri fra i suoi estremi, mentre la durata cambia con la cronologia. Una finestra di più o meno cento anni mantiene invece un criterio temporale e può contenere un numero diverso di campioni nelle varie porzioni della successione. La media dei punti entro la finestra attribuisce uguale peso ai campioni, mentre una media sul tempo richiederebbe pesi coerenti con gli intervalli rappresentati. Il metodo di smoothing deve quindi specificare ampiezza, unità e regola di ponderazione.\par\vspace{.28cm}{\centering\fontsize{13}{17}\selectfont\(\displaystyle\begin{gathered}\mathcal W_t=\{i:|t_i-t_0|\le100\ \mathrm{anni}\}\end{gathered}\)\par}\end{minipage}\end{lrbox}
\begin{lrbox}{\SideBox}\begin{minipage}[t]{7.45cm}\raggedright\sloppy\color{Ink}\fontsize{12.2}{16.1}\selectfont \input{tikz/M13/D07.tex}\par\vspace{.16cm}{\fontsize{9.5}{12.6}\selectfont $t_0$: età centrale; $t_i$: età assegnata al campione i. La finestra temporale include tutti i campioni entro 100 anni dal centro.}\end{minipage}\end{lrbox}
\begin{lrbox}{\RecallBox}\begin{minipage}[t]{16.45cm}\raggedright\sloppy\color{Ink}\fontsize{10.1}{13.2}\selectfont L’unità della finestra e i pesi stabiliscono quale variazione venga attenuata.\end{minipage}\end{lrbox}
\MeasureBox{19-TitleBox}{\TitleBox}
\MeasureBox{19-BodyBox}{\BodyBox}
\MeasureBox{19-SideBox}{\SideBox}
\MeasureBox{19-RecallBox}{\RecallBox}
\begin{frame}[plain]\hypertarget{slide19}{}
\bookmark[level=1,page=19]{Finestre per indice e finestre temporali}
\begin{tikzpicture}[remember picture,overlay]\node[anchor=north east,inner sep=0] at ([xshift=19.05cm,yshift=-0.16cm]current page.north west){\color{Classical}\fontsize{8}{10}\selectfont\hyperlink{slide2}{Indice}};
\node[anchor=north west,inner sep=0] at ([xshift=0.95cm,yshift=-0.62cm]current page.north west){\usebox{\TitleBox}};
\draw[Classical,line width=1.3pt] ([xshift=.95cm,yshift=-1.68cm]current page.north west)--([xshift=2.25cm,yshift=-1.68cm]current page.north west);\node[anchor=north west,inner sep=0] at ([xshift=0.95cm,yshift=-2.12cm]current page.north west){\usebox{\BodyBox}};
\node[anchor=north west,inner sep=0] at ([xshift=11.6cm,yshift=-2.1cm]current page.north west){\color{Quantum}\bfseries\fontsize{11.5}{14}\selectfont Lettura del diagramma};
\node[anchor=north west,inner sep=0] at ([xshift=11.6cm,yshift=-2.75cm]current page.north west){\usebox{\SideBox}};
\draw[Rule,line width=.5pt] ([xshift=.95cm,yshift=-10.43cm]current page.north west)--([xshift=19.05cm,yshift=-10.43cm]current page.north west);\node[anchor=north west,inner sep=0] at ([xshift=0.95cm,yshift=-10.66cm]current page.north west){\color{Classical}\bfseries\fontsize{9.4}{12}\selectfont Richiami};
\node[anchor=north west,inner sep=0] at ([xshift=2.6cm,yshift=-10.64cm]current page.north west){\usebox{\RecallBox}};
\draw[Rule,line width=.35pt] ([xshift=.95cm,yshift=-11.92cm]current page.north west)--([xshift=19.05cm,yshift=-11.92cm]current page.north west);\node[anchor=north west,inner sep=0] at ([xshift=0.95cm,yshift=-12.06cm]current page.north west){\color{Muted}\fontsize{7.7}{9.3}\selectfont M13 · Cronologie e archivi sedimentari};
\node[anchor=north east,inner sep=0] at ([xshift=17.65cm,yshift=-12.06cm]current page.north west){\color{Classical}\fontsize{7.3}{9.1}\selectfont };
\node[anchor=north east,inner sep=0] at ([xshift=19.05cm,yshift=-12.06cm]current page.north west){\color{Muted}\fontsize{7.7}{9.3}\selectfont 19/30};
\end{tikzpicture}\end{frame}
\begin{lrbox}{\TitleBox}\begin{minipage}[t]{18.2cm}\raggedright\sloppy\color{Ink}\fontsize{20.7}{24}\selectfont \bfseries\color{Navy}Interpolazione temporale e informazione osservata\end{minipage}\end{lrbox}
\begin{lrbox}{\BodyBox}\begin{minipage}[t]{10.25cm}\raggedright\sloppy\color{Ink}\fontsize{12.4}{16.4}\selectfont Ricostruire il proxy ogni dieci anni produce 221 righe, tutte derivate dalle ottantuno misure originarie. La griglia regolare può essere utile per confronti grafici o per procedure che richiedano un asse uniforme, purché si riconoscano la dipendenza introdotta e le ipotesi dell’interpolazione. Il numero delle righe ricostruite non fornisce nuovi gradi di libertà osservativi. Se l’analisi successiva riguarda incertezza o significatività, deve conservare la provenienza dai campioni originali e la variabilità della trasformazione cronologica.\par\vspace{.28cm}{\centering\fontsize{13}{17}\selectfont\(\displaystyle\begin{gathered}\widetilde y(q)=\sum_i a_i(q)y_i\end{gathered}\)\par}\end{minipage}\end{lrbox}
\begin{lrbox}{\SideBox}\begin{minipage}[t]{7.45cm}\raggedright\sloppy\color{Ink}\fontsize{12.2}{16.1}\selectfont q è l’età richiesta; $y_i$ il proxy osservato; $a_i(q)$ il peso del campione i. Nell’interpolazione lineare interna intervengono due vicini, con pesi non negativi che sommano a uno. Più punti ricostruiti possono condividere campioni e ancore cronologiche, introducendo dipendenza su entrambi gli assi.\end{minipage}\end{lrbox}
\begin{lrbox}{\RecallBox}\begin{minipage}[t]{16.45cm}\raggedright\sloppy\color{Ink}\fontsize{10.1}{13.2}\selectfont La regolarizzazione della griglia temporale modifica la rappresentazione e il modello di dipendenza.\end{minipage}\end{lrbox}
\MeasureBox{20-TitleBox}{\TitleBox}
\MeasureBox{20-BodyBox}{\BodyBox}
\MeasureBox{20-SideBox}{\SideBox}
\MeasureBox{20-RecallBox}{\RecallBox}
\begin{frame}[plain]\hypertarget{slide20}{}
\bookmark[level=1,page=20]{Interpolazione temporale e informazione osservata}
\begin{tikzpicture}[remember picture,overlay]\node[anchor=north east,inner sep=0] at ([xshift=19.05cm,yshift=-0.16cm]current page.north west){\color{Classical}\fontsize{8}{10}\selectfont\hyperlink{slide2}{Indice}};
\node[anchor=north west,inner sep=0] at ([xshift=0.95cm,yshift=-0.62cm]current page.north west){\usebox{\TitleBox}};
\draw[Classical,line width=1.3pt] ([xshift=.95cm,yshift=-1.68cm]current page.north west)--([xshift=2.25cm,yshift=-1.68cm]current page.north west);\node[anchor=north west,inner sep=0] at ([xshift=0.95cm,yshift=-2.12cm]current page.north west){\usebox{\BodyBox}};
\node[anchor=north west,inner sep=0] at ([xshift=11.6cm,yshift=-2.1cm]current page.north west){\color{Quantum}\bfseries\fontsize{11.5}{14}\selectfont Sviluppo e interpretazione};
\node[anchor=north west,inner sep=0] at ([xshift=11.6cm,yshift=-2.75cm]current page.north west){\usebox{\SideBox}};
\draw[Rule,line width=.5pt] ([xshift=.95cm,yshift=-10.43cm]current page.north west)--([xshift=19.05cm,yshift=-10.43cm]current page.north west);\node[anchor=north west,inner sep=0] at ([xshift=0.95cm,yshift=-10.66cm]current page.north west){\color{Classical}\bfseries\fontsize{9.4}{12}\selectfont Richiami};
\node[anchor=north west,inner sep=0] at ([xshift=2.6cm,yshift=-10.64cm]current page.north west){\usebox{\RecallBox}};
\draw[Rule,line width=.35pt] ([xshift=.95cm,yshift=-11.92cm]current page.north west)--([xshift=19.05cm,yshift=-11.92cm]current page.north west);\node[anchor=north west,inner sep=0] at ([xshift=0.95cm,yshift=-12.06cm]current page.north west){\color{Muted}\fontsize{7.7}{9.3}\selectfont M13 · Cronologie e archivi sedimentari};
\node[anchor=north east,inner sep=0] at ([xshift=17.65cm,yshift=-12.06cm]current page.north west){\color{Classical}\fontsize{7.3}{9.1}\selectfont };
\node[anchor=north east,inner sep=0] at ([xshift=19.05cm,yshift=-12.06cm]current page.north west){\color{Muted}\fontsize{7.7}{9.3}\selectfont 20/30};
\end{tikzpicture}\end{frame}
\begin{lrbox}{\TitleBox}\begin{minipage}[t]{18.2cm}\raggedright\sloppy\color{Ink}\fontsize{20.7}{24}\selectfont \bfseries\color{Navy}Hiatus come discontinuità della cronologia\end{minipage}\end{lrbox}
\begin{lrbox}{\BodyBox}\begin{minipage}[t]{10.25cm}\raggedright\sloppy\color{Ink}\fontsize{12.4}{16.4}\selectfont Lo scenario di hiatus aggiunge quattrocento anni alle profondità superiori a ottanta centimetri, mantenendo separata la cronologia di base. Il salto rappresenta un intervallo temporale privo di sedimentazione conservata nel modello, per il quale servirebbero evidenze stratigrafiche o cronologiche indipendenti. La figura distingue i due rami ai lati della discontinuità, evitando che una linea obliqua suggerisca accumulo attraverso l’intervallo mancante. La scelta della profondità e della durata del salto appartiene allo scenario e non viene inferita automaticamente dal proxy.\par\vspace{.28cm}{\centering\fontsize{13}{17}\selectfont\(\displaystyle\begin{gathered}t_H(z)=t(z)+400\,\mathbf1\{z>80\}\end{gathered}\)\par}\end{minipage}\end{lrbox}
\begin{lrbox}{\SideBox}\begin{minipage}[t]{7.45cm}\raggedright\sloppy\color{Ink}\fontsize{12.2}{16.1}\selectfont \includegraphics[width=7.45cm]{figure/M13/F07.pdf}\par\vspace{.16cm}{\fontsize{9.5}{12.6}\selectfont Verde: cronologia di base. Viola: scenario con hiatus; il tratteggio indica il salto a 80 centimetri.}\end{minipage}\end{lrbox}
\begin{lrbox}{\RecallBox}\begin{minipage}[t]{16.45cm}\raggedright\sloppy\color{Ink}\fontsize{10.1}{13.2}\selectfont Una lacuna deve essere sostenuta da evidenze indipendenti e rappresentata come discontinuità esplicita.\end{minipage}\end{lrbox}
\MeasureBox{21-TitleBox}{\TitleBox}
\MeasureBox{21-BodyBox}{\BodyBox}
\MeasureBox{21-SideBox}{\SideBox}
\MeasureBox{21-RecallBox}{\RecallBox}
\begin{frame}[plain]\hypertarget{slide21}{}
\bookmark[level=0,page=21]{Discontinuità e sensibilità del modello}
\bookmark[level=1,page=21]{Hiatus come discontinuità della cronologia}
\begin{tikzpicture}[remember picture,overlay]\node[anchor=north east,inner sep=0] at ([xshift=19.05cm,yshift=-0.16cm]current page.north west){\color{Classical}\fontsize{8}{10}\selectfont\hyperlink{slide2}{Indice}};
\node[anchor=north west,inner sep=0] at ([xshift=0.95cm,yshift=-0.62cm]current page.north west){\usebox{\TitleBox}};
\draw[Classical,line width=1.3pt] ([xshift=.95cm,yshift=-1.68cm]current page.north west)--([xshift=2.25cm,yshift=-1.68cm]current page.north west);\node[anchor=north west,inner sep=0] at ([xshift=0.95cm,yshift=-2.12cm]current page.north west){\usebox{\BodyBox}};
\node[anchor=north west,inner sep=0] at ([xshift=11.6cm,yshift=-2.1cm]current page.north west){\color{Quantum}\bfseries\fontsize{11.5}{14}\selectfont Evidenza grafica};
\node[anchor=north west,inner sep=0] at ([xshift=11.6cm,yshift=-2.75cm]current page.north west){\usebox{\SideBox}};
\draw[Rule,line width=.5pt] ([xshift=.95cm,yshift=-10.43cm]current page.north west)--([xshift=19.05cm,yshift=-10.43cm]current page.north west);\node[anchor=north west,inner sep=0] at ([xshift=0.95cm,yshift=-10.66cm]current page.north west){\color{Classical}\bfseries\fontsize{9.4}{12}\selectfont Richiami};
\node[anchor=north west,inner sep=0] at ([xshift=2.6cm,yshift=-10.64cm]current page.north west){\usebox{\RecallBox}};
\draw[Rule,line width=.35pt] ([xshift=.95cm,yshift=-11.92cm]current page.north west)--([xshift=19.05cm,yshift=-11.92cm]current page.north west);\node[anchor=north west,inner sep=0] at ([xshift=0.95cm,yshift=-12.06cm]current page.north west){\color{Muted}\fontsize{7.7}{9.3}\selectfont M13 · Cronologie e archivi sedimentari};
\node[anchor=north east,inner sep=0] at ([xshift=17.65cm,yshift=-12.06cm]current page.north west){\color{Classical}\fontsize{7.3}{9.1}\selectfont };
\node[anchor=north east,inner sep=0] at ([xshift=19.05cm,yshift=-12.06cm]current page.north west){\color{Muted}\fontsize{7.7}{9.3}\selectfont 21/30};
\end{tikzpicture}\end{frame}
\begin{lrbox}{\TitleBox}\begin{minipage}[t]{18.2cm}\raggedright\sloppy\color{Ink}\fontsize{20.7}{24}\selectfont \bfseries\color{Navy}Inversione cronologica e diagnosi dell’incongruenza\end{minipage}\end{lrbox}
\begin{lrbox}{\BodyBox}\begin{minipage}[t]{10.25cm}\raggedright\sloppy\color{Ink}\fontsize{12.4}{16.4}\selectfont Sostituendo in una copia l’età a ottanta centimetri con trecento anni, si introduce un’inversione fra cinquanta e ottanta centimetri. Il controllo individua il segmento nel quale l’età diminuisce con la profondità, mentre la spiegazione può richiedere un esame di errore analitico, rimaneggiamento, disturbo o associazione errata. Eliminare l’ancora soltanto perché ostacola la monotonicità trasferirebbe nel risultato una scelta non documentata. Eventuali esclusioni devono conservare il dato originario, la motivazione e un confronto di sensibilità.\par\vspace{.28cm}{\centering\fontsize{13}{17}\selectfont\(\displaystyle\begin{gathered}\Delta t_i\le0\quad\Rightarrow\quad\text{verifica del segmento }i\end{gathered}\)\par}\end{minipage}\end{lrbox}
\begin{lrbox}{\SideBox}\begin{minipage}[t]{7.45cm}\raggedright\sloppy\color{Ink}\fontsize{12.2}{16.1}\selectfont \includegraphics[width=7.45cm]{figure/M13/F08.pdf}\par\vspace{.16cm}{\fontsize{9.5}{12.6}\selectfont Una sola ancora modificata produce un tratto cronologicamente invertito. La figura conserva l’ordine delle profondità.}\end{minipage}\end{lrbox}
\begin{lrbox}{\RecallBox}\begin{minipage}[t]{16.45cm}\raggedright\sloppy\color{Ink}\fontsize{10.1}{13.2}\selectfont Il controllo riconosce l’incongruenza, mentre la diagnosi richiede evidenze ulteriori.\end{minipage}\end{lrbox}
\MeasureBox{22-TitleBox}{\TitleBox}
\MeasureBox{22-BodyBox}{\BodyBox}
\MeasureBox{22-SideBox}{\SideBox}
\MeasureBox{22-RecallBox}{\RecallBox}
\begin{frame}[plain]\hypertarget{slide22}{}
\bookmark[level=1,page=22]{Inversione cronologica e diagnosi dell’incongruenza}
\begin{tikzpicture}[remember picture,overlay]\node[anchor=north east,inner sep=0] at ([xshift=19.05cm,yshift=-0.16cm]current page.north west){\color{Classical}\fontsize{8}{10}\selectfont\hyperlink{slide2}{Indice}};
\node[anchor=north west,inner sep=0] at ([xshift=0.95cm,yshift=-0.62cm]current page.north west){\usebox{\TitleBox}};
\draw[Classical,line width=1.3pt] ([xshift=.95cm,yshift=-1.68cm]current page.north west)--([xshift=2.25cm,yshift=-1.68cm]current page.north west);\node[anchor=north west,inner sep=0] at ([xshift=0.95cm,yshift=-2.12cm]current page.north west){\usebox{\BodyBox}};
\node[anchor=north west,inner sep=0] at ([xshift=11.6cm,yshift=-2.1cm]current page.north west){\color{Quantum}\bfseries\fontsize{11.5}{14}\selectfont Evidenza grafica};
\node[anchor=north west,inner sep=0] at ([xshift=11.6cm,yshift=-2.75cm]current page.north west){\usebox{\SideBox}};
\draw[Rule,line width=.5pt] ([xshift=.95cm,yshift=-10.43cm]current page.north west)--([xshift=19.05cm,yshift=-10.43cm]current page.north west);\node[anchor=north west,inner sep=0] at ([xshift=0.95cm,yshift=-10.66cm]current page.north west){\color{Classical}\bfseries\fontsize{9.4}{12}\selectfont Richiami};
\node[anchor=north west,inner sep=0] at ([xshift=2.6cm,yshift=-10.64cm]current page.north west){\usebox{\RecallBox}};
\draw[Rule,line width=.35pt] ([xshift=.95cm,yshift=-11.92cm]current page.north west)--([xshift=19.05cm,yshift=-11.92cm]current page.north west);\node[anchor=north west,inner sep=0] at ([xshift=0.95cm,yshift=-12.06cm]current page.north west){\color{Muted}\fontsize{7.7}{9.3}\selectfont M13 · Cronologie e archivi sedimentari};
\node[anchor=north east,inner sep=0] at ([xshift=17.65cm,yshift=-12.06cm]current page.north west){\color{Classical}\fontsize{7.3}{9.1}\selectfont };
\node[anchor=north east,inner sep=0] at ([xshift=19.05cm,yshift=-12.06cm]current page.north west){\color{Muted}\fontsize{7.7}{9.3}\selectfont 22/30};
\end{tikzpicture}\end{frame}
\begin{lrbox}{\TitleBox}\begin{minipage}[t]{18.2cm}\raggedright\sloppy\color{Ink}\fontsize{20.7}{24}\selectfont \bfseries\color{Navy}Retta globale e variabilità dei tassi locali\end{minipage}\end{lrbox}
\begin{lrbox}{\BodyBox}\begin{minipage}[t]{10.25cm}\raggedright\sloppy\color{Ink}\fontsize{12.4}{16.4}\selectfont Una regressione lineare età–profondità assume un unico rapporto medio fra le due coordinate, mentre l’interpolazione a tratti consente cambiamenti di pendenza in corrispondenza delle ancore. La retta può avere residui anche quando la successione è monotona, poiché i tassi apparenti variano. L’interpolatore attraversa invece tutte le ancore per costruzione, cosicché il suo errore di adattamento nullo non dimostra migliore previsione. La scelta deve essere confrontata con il processo deposizionale, con la qualità delle ancore e con controlli esterni o di sensibilità.\par\vspace{.28cm}{\centering\fontsize{13}{17}\selectfont\(\displaystyle\begin{gathered}t(z)=\beta_0+\beta_1z+\varepsilon\end{gathered}\)\par}\end{minipage}\end{lrbox}
\begin{lrbox}{\SideBox}\begin{minipage}[t]{7.45cm}\raggedright\sloppy\color{Ink}\fontsize{12.2}{16.1}\selectfont \includegraphics[width=7.45cm]{figure/M13/F09.pdf}\par\vspace{.16cm}{\fontsize{9.5}{12.6}\selectfont $\beta_0$: intercetta; $\beta_1$: anni per cm; $\varepsilon$: scarto. Verde: interpolazione a tratti; viola: retta sulle sei ancore.}\end{minipage}\end{lrbox}
\begin{lrbox}{\RecallBox}\begin{minipage}[t]{16.45cm}\raggedright\sloppy\color{Ink}\fontsize{10.1}{13.2}\selectfont La flessibilità della curva deve essere giustificata dal processo e dalla capacità di previsione.\end{minipage}\end{lrbox}
\MeasureBox{23-TitleBox}{\TitleBox}
\MeasureBox{23-BodyBox}{\BodyBox}
\MeasureBox{23-SideBox}{\SideBox}
\MeasureBox{23-RecallBox}{\RecallBox}
\begin{frame}[plain]\hypertarget{slide23}{}
\bookmark[level=1,page=23]{Retta globale e variabilità dei tassi locali}
\begin{tikzpicture}[remember picture,overlay]\node[anchor=north east,inner sep=0] at ([xshift=19.05cm,yshift=-0.16cm]current page.north west){\color{Classical}\fontsize{8}{10}\selectfont\hyperlink{slide2}{Indice}};
\node[anchor=north west,inner sep=0] at ([xshift=0.95cm,yshift=-0.62cm]current page.north west){\usebox{\TitleBox}};
\draw[Classical,line width=1.3pt] ([xshift=.95cm,yshift=-1.68cm]current page.north west)--([xshift=2.25cm,yshift=-1.68cm]current page.north west);\node[anchor=north west,inner sep=0] at ([xshift=0.95cm,yshift=-2.12cm]current page.north west){\usebox{\BodyBox}};
\node[anchor=north west,inner sep=0] at ([xshift=11.6cm,yshift=-2.1cm]current page.north west){\color{Quantum}\bfseries\fontsize{11.5}{14}\selectfont Evidenza grafica};
\node[anchor=north west,inner sep=0] at ([xshift=11.6cm,yshift=-2.75cm]current page.north west){\usebox{\SideBox}};
\draw[Rule,line width=.5pt] ([xshift=.95cm,yshift=-10.43cm]current page.north west)--([xshift=19.05cm,yshift=-10.43cm]current page.north west);\node[anchor=north west,inner sep=0] at ([xshift=0.95cm,yshift=-10.66cm]current page.north west){\color{Classical}\bfseries\fontsize{9.4}{12}\selectfont Richiami};
\node[anchor=north west,inner sep=0] at ([xshift=2.6cm,yshift=-10.64cm]current page.north west){\usebox{\RecallBox}};
\draw[Rule,line width=.35pt] ([xshift=.95cm,yshift=-11.92cm]current page.north west)--([xshift=19.05cm,yshift=-11.92cm]current page.north west);\node[anchor=north west,inner sep=0] at ([xshift=0.95cm,yshift=-12.06cm]current page.north west){\color{Muted}\fontsize{7.7}{9.3}\selectfont M13 · Cronologie e archivi sedimentari};
\node[anchor=north east,inner sep=0] at ([xshift=17.65cm,yshift=-12.06cm]current page.north west){\color{Classical}\fontsize{7.3}{9.1}\selectfont };
\node[anchor=north east,inner sep=0] at ([xshift=19.05cm,yshift=-12.06cm]current page.north west){\color{Muted}\fontsize{7.7}{9.3}\selectfont 23/30};
\end{tikzpicture}\end{frame}
\begin{lrbox}{\TitleBox}\begin{minipage}[t]{18.2cm}\raggedright\sloppy\color{Ink}\fontsize{20.7}{24}\selectfont \bfseries\color{Navy}Esclusione di un’ancora interna e sensibilità locale\end{minipage}\end{lrbox}
\begin{lrbox}{\BodyBox}\begin{minipage}[t]{10.25cm}\raggedright\sloppy\color{Ink}\fontsize{12.4}{16.4}\selectfont Si elimina un’ancora interna alla volta e si stima la sua età mediante le ancore rimaste. Gli estremi vengono conservati, perché la loro esclusione richiederebbe estrapolazione. Gli scarti risultanti misurano quanto la ricostruzione locale dipenda dall’ancora omessa, senza riprodurre l’acquisizione di nuove datazioni indipendenti. Uno scarto ampio può segnalare una variazione reale del tasso, oltre a un possibile problema della misura. La verifica deve quindi orientare l’approfondimento e non diventare una regola automatica di eliminazione.\par\vspace{.28cm}{\centering\fontsize{13}{17}\selectfont\(\displaystyle\begin{gathered}e_{(-i)}=\widehat t_{(-i)}(z_i)-t_i,\qquad i=2,\ldots,5\end{gathered}\)\par}\end{minipage}\end{lrbox}
\begin{lrbox}{\SideBox}\begin{minipage}[t]{7.45cm}\raggedright\sloppy\color{Ink}\fontsize{12.2}{16.1}\selectfont \includegraphics[width=7.45cm]{figure/M13/F10.pdf}\par\vspace{.16cm}{\fontsize{9.5}{12.6}\selectfont Scarti delle quattro ancore interne dopo l’esclusione individuale. Il segno è previsto meno osservato.}\end{minipage}\end{lrbox}
\begin{lrbox}{\RecallBox}\begin{minipage}[t]{16.45cm}\raggedright\sloppy\color{Ink}\fontsize{10.1}{13.2}\selectfont La sensibilità a un’ancora va interpretata nel contesto della successione e del modello.\end{minipage}\end{lrbox}
\MeasureBox{24-TitleBox}{\TitleBox}
\MeasureBox{24-BodyBox}{\BodyBox}
\MeasureBox{24-SideBox}{\SideBox}
\MeasureBox{24-RecallBox}{\RecallBox}
\begin{frame}[plain]\hypertarget{slide24}{}
\bookmark[level=1,page=24]{Esclusione di un’ancora interna e sensibilità locale}
\begin{tikzpicture}[remember picture,overlay]\node[anchor=north east,inner sep=0] at ([xshift=19.05cm,yshift=-0.16cm]current page.north west){\color{Classical}\fontsize{8}{10}\selectfont\hyperlink{slide2}{Indice}};
\node[anchor=north west,inner sep=0] at ([xshift=0.95cm,yshift=-0.62cm]current page.north west){\usebox{\TitleBox}};
\draw[Classical,line width=1.3pt] ([xshift=.95cm,yshift=-1.68cm]current page.north west)--([xshift=2.25cm,yshift=-1.68cm]current page.north west);\node[anchor=north west,inner sep=0] at ([xshift=0.95cm,yshift=-2.12cm]current page.north west){\usebox{\BodyBox}};
\node[anchor=north west,inner sep=0] at ([xshift=11.6cm,yshift=-2.1cm]current page.north west){\color{Quantum}\bfseries\fontsize{11.5}{14}\selectfont Evidenza grafica};
\node[anchor=north west,inner sep=0] at ([xshift=11.6cm,yshift=-2.75cm]current page.north west){\usebox{\SideBox}};
\draw[Rule,line width=.5pt] ([xshift=.95cm,yshift=-10.43cm]current page.north west)--([xshift=19.05cm,yshift=-10.43cm]current page.north west);\node[anchor=north west,inner sep=0] at ([xshift=0.95cm,yshift=-10.66cm]current page.north west){\color{Classical}\bfseries\fontsize{9.4}{12}\selectfont Richiami};
\node[anchor=north west,inner sep=0] at ([xshift=2.6cm,yshift=-10.64cm]current page.north west){\usebox{\RecallBox}};
\draw[Rule,line width=.35pt] ([xshift=.95cm,yshift=-11.92cm]current page.north west)--([xshift=19.05cm,yshift=-11.92cm]current page.north west);\node[anchor=north west,inner sep=0] at ([xshift=0.95cm,yshift=-12.06cm]current page.north west){\color{Muted}\fontsize{7.7}{9.3}\selectfont M13 · Cronologie e archivi sedimentari};
\node[anchor=north east,inner sep=0] at ([xshift=17.65cm,yshift=-12.06cm]current page.north west){\color{Classical}\fontsize{7.3}{9.1}\selectfont };
\node[anchor=north east,inner sep=0] at ([xshift=19.05cm,yshift=-12.06cm]current page.north west){\color{Muted}\fontsize{7.7}{9.3}\selectfont 24/30};
\end{tikzpicture}\end{frame}
\begin{lrbox}{\TitleBox}\begin{minipage}[t]{18.2cm}\raggedright\sloppy\color{Ink}\fontsize{20.7}{24}\selectfont \bfseries\color{Navy}Derivata del proxy e cambiamento di scala\end{minipage}\end{lrbox}
\begin{lrbox}{\BodyBox}\begin{minipage}[t]{10.25cm}\raggedright\sloppy\color{Ink}\fontsize{12.4}{16.4}\selectfont La variazione del proxy per centimetro descrive il cambiamento lungo lo spessore, mentre la variazione per anno dipende anche dalla cronologia. Dove le funzioni sono derivabili e il tasso è positivo, la regola della catena collega le due derivate tramite il tasso locale di accumulo. Una conversione basata su un unico tasso medio è adeguata soltanto se la trasformazione dell’asse è coerente con tale approssimazione. Due regressioni globali sui diversi assi possono quindi restituire pendenze che non si ottengono mediante una semplice riscalatura costante.\par\vspace{.28cm}{\centering\fontsize{13}{17}\selectfont\(\displaystyle\begin{gathered}\frac{\mathrm dy}{\mathrm dt}=\frac{\mathrm dy}{\mathrm dz}\frac{\mathrm dz}{\mathrm dt}=r(z)\frac{\mathrm dy}{\mathrm dz}\end{gathered}\)\par}\end{minipage}\end{lrbox}
\begin{lrbox}{\SideBox}\begin{minipage}[t]{7.45cm}\raggedright\sloppy\color{Ink}\fontsize{12.2}{16.1}\selectfont \input{tikz/M13/D08.tex}\par\vspace{.16cm}{\fontsize{9.5}{12.6}\selectfont y: proxy in u.a.; z: profondità in cm; t: età in anni; r(z): tasso in cm/anno. La derivata temporale è in u.a./anno.}\end{minipage}\end{lrbox}
\begin{lrbox}{\RecallBox}\begin{minipage}[t]{16.45cm}\raggedright\sloppy\color{Ink}\fontsize{10.1}{13.2}\selectfont La trasformazione dell’asse cambia l’unità e l’interpretazione della variazione.\end{minipage}\end{lrbox}
\MeasureBox{25-TitleBox}{\TitleBox}
\MeasureBox{25-BodyBox}{\BodyBox}
\MeasureBox{25-SideBox}{\SideBox}
\MeasureBox{25-RecallBox}{\RecallBox}
\begin{frame}[plain]\hypertarget{slide25}{}
\bookmark[level=0,page=25]{Pendenze, durate ed eventi}
\bookmark[level=1,page=25]{Derivata del proxy e cambiamento di scala}
\begin{tikzpicture}[remember picture,overlay]\node[anchor=north east,inner sep=0] at ([xshift=19.05cm,yshift=-0.16cm]current page.north west){\color{Classical}\fontsize{8}{10}\selectfont\hyperlink{slide2}{Indice}};
\node[anchor=north west,inner sep=0] at ([xshift=0.95cm,yshift=-0.62cm]current page.north west){\usebox{\TitleBox}};
\draw[Classical,line width=1.3pt] ([xshift=.95cm,yshift=-1.68cm]current page.north west)--([xshift=2.25cm,yshift=-1.68cm]current page.north west);\node[anchor=north west,inner sep=0] at ([xshift=0.95cm,yshift=-2.12cm]current page.north west){\usebox{\BodyBox}};
\node[anchor=north west,inner sep=0] at ([xshift=11.6cm,yshift=-2.1cm]current page.north west){\color{Quantum}\bfseries\fontsize{11.5}{14}\selectfont Lettura del diagramma};
\node[anchor=north west,inner sep=0] at ([xshift=11.6cm,yshift=-2.75cm]current page.north west){\usebox{\SideBox}};
\draw[Rule,line width=.5pt] ([xshift=.95cm,yshift=-10.43cm]current page.north west)--([xshift=19.05cm,yshift=-10.43cm]current page.north west);\node[anchor=north west,inner sep=0] at ([xshift=0.95cm,yshift=-10.66cm]current page.north west){\color{Classical}\bfseries\fontsize{9.4}{12}\selectfont Richiami};
\node[anchor=north west,inner sep=0] at ([xshift=2.6cm,yshift=-10.64cm]current page.north west){\usebox{\RecallBox}};
\draw[Rule,line width=.35pt] ([xshift=.95cm,yshift=-11.92cm]current page.north west)--([xshift=19.05cm,yshift=-11.92cm]current page.north west);\node[anchor=north west,inner sep=0] at ([xshift=0.95cm,yshift=-12.06cm]current page.north west){\color{Muted}\fontsize{7.7}{9.3}\selectfont M13 · Cronologie e archivi sedimentari};
\node[anchor=north east,inner sep=0] at ([xshift=17.65cm,yshift=-12.06cm]current page.north west){\color{Classical}\fontsize{7.3}{9.1}\selectfont };
\node[anchor=north east,inner sep=0] at ([xshift=19.05cm,yshift=-12.06cm]current page.north west){\color{Muted}\fontsize{7.7}{9.3}\selectfont 25/30};
\end{tikzpicture}\end{frame}
\begin{lrbox}{\TitleBox}\begin{minipage}[t]{18.2cm}\raggedright\sloppy\color{Ink}\fontsize{20.7}{24}\selectfont \bfseries\color{Navy}Durata di un intervallo e covarianza condivisa\end{minipage}\end{lrbox}
\begin{lrbox}{\BodyBox}\begin{minipage}[t]{10.25cm}\raggedright\sloppy\color{Ink}\fontsize{12.4}{16.4}\selectfont La durata fra sessanta e cento centimetri si calcola sottraendo le due età all’interno dello stesso scenario. Le età condividono l’ancora a ottanta centimetri, per cui la loro covarianza entra nella varianza della differenza. Sottrarre quantili marginali costruirebbe una quantità diversa dai quantili della durata e perderebbe l’abbinamento fra scenari. Prima del condizionamento, la deviazione standard analitica è circa 48,94 anni; il codice verifica l’equivalenza delle forme matriciale ed espansa.\par\vspace{.28cm}{\centering\fontsize{13}{17}\selectfont\(\displaystyle\begin{gathered}D=T_2-T_1,\quad T_1=T(60),\ T_2=T(100)\\ \operatorname{Var}(D)=\sigma_1^2+\sigma_2^2-2\sigma_{12}\end{gathered}\)\par}\end{minipage}\end{lrbox}
\begin{lrbox}{\SideBox}\begin{minipage}[t]{7.45cm}\raggedright\sloppy\color{Ink}\fontsize{12.2}{16.1}\selectfont Nella formula, i termini quadratici indicano le varianze delle due età; il termine con indici 12 indica la loro covarianza. I pesi delle due profondità sono differenti: a 60 cm si combinano le ancore di 50 e 80 cm; a 100 cm si combinano quelle di 80 e 120 cm. La perturbazione dell’ancora condivisa entra in entrambe le età e viene parzialmente compensata nella differenza.\end{minipage}\end{lrbox}
\begin{lrbox}{\RecallBox}\begin{minipage}[t]{16.45cm}\raggedright\sloppy\color{Ink}\fontsize{10.1}{13.2}\selectfont Le quantità derivate da una cronologia devono conservare la struttura congiunta delle età.\end{minipage}\end{lrbox}
\MeasureBox{26-TitleBox}{\TitleBox}
\MeasureBox{26-BodyBox}{\BodyBox}
\MeasureBox{26-SideBox}{\SideBox}
\MeasureBox{26-RecallBox}{\RecallBox}
\begin{frame}[plain]\hypertarget{slide26}{}
\bookmark[level=1,page=26]{Durata di un intervallo e covarianza condivisa}
\begin{tikzpicture}[remember picture,overlay]\node[anchor=north east,inner sep=0] at ([xshift=19.05cm,yshift=-0.16cm]current page.north west){\color{Classical}\fontsize{8}{10}\selectfont\hyperlink{slide2}{Indice}};
\node[anchor=north west,inner sep=0] at ([xshift=0.95cm,yshift=-0.62cm]current page.north west){\usebox{\TitleBox}};
\draw[Classical,line width=1.3pt] ([xshift=.95cm,yshift=-1.68cm]current page.north west)--([xshift=2.25cm,yshift=-1.68cm]current page.north west);\node[anchor=north west,inner sep=0] at ([xshift=0.95cm,yshift=-2.12cm]current page.north west){\usebox{\BodyBox}};
\node[anchor=north west,inner sep=0] at ([xshift=11.6cm,yshift=-2.1cm]current page.north west){\color{Quantum}\bfseries\fontsize{11.5}{14}\selectfont Sviluppo e interpretazione};
\node[anchor=north west,inner sep=0] at ([xshift=11.6cm,yshift=-2.75cm]current page.north west){\usebox{\SideBox}};
\draw[Rule,line width=.5pt] ([xshift=.95cm,yshift=-10.43cm]current page.north west)--([xshift=19.05cm,yshift=-10.43cm]current page.north west);\node[anchor=north west,inner sep=0] at ([xshift=0.95cm,yshift=-10.66cm]current page.north west){\color{Classical}\bfseries\fontsize{9.4}{12}\selectfont Richiami};
\node[anchor=north west,inner sep=0] at ([xshift=2.6cm,yshift=-10.64cm]current page.north west){\usebox{\RecallBox}};
\draw[Rule,line width=.35pt] ([xshift=.95cm,yshift=-11.92cm]current page.north west)--([xshift=19.05cm,yshift=-11.92cm]current page.north west);\node[anchor=north west,inner sep=0] at ([xshift=0.95cm,yshift=-12.06cm]current page.north west){\color{Muted}\fontsize{7.7}{9.3}\selectfont M13 · Cronologie e archivi sedimentari};
\node[anchor=north east,inner sep=0] at ([xshift=17.65cm,yshift=-12.06cm]current page.north west){\color{Classical}\fontsize{7.3}{9.1}\selectfont };
\node[anchor=north east,inner sep=0] at ([xshift=19.05cm,yshift=-12.06cm]current page.north west){\color{Muted}\fontsize{7.7}{9.3}\selectfont 26/30};
\end{tikzpicture}\end{frame}
\begin{lrbox}{\TitleBox}\begin{minipage}[t]{18.2cm}\raggedright\sloppy\color{Ink}\fontsize{20.7}{24}\selectfont \bfseries\color{Navy}Distribuzione degli scenari di durata\end{minipage}\end{lrbox}
\begin{lrbox}{\BodyBox}\begin{minipage}[t]{10.25cm}\raggedright\sloppy\color{Ink}\fontsize{12.4}{16.4}\selectfont La cronologia di base attribuisce 530 anni all’intervallo fra sessanta e cento centimetri. Negli scenari inclusi, la mediana è circa 529,7 anni e i quantili al 2,5\% e 97,5\% sono circa 433,0 e 630,5 anni. Il frammento ricostruisce la differenza per ciascun vettore di ancore, preservando la dipendenza richiesta. La distribuzione descrive soltanto le perturbazioni adottate e la regola lineare, mentre un hiatus entro l’intervallo o un diverso modello di accumulo modificherebbero il significato e il valore della durata.\end{minipage}\end{lrbox}
\begin{lrbox}{\SideBox}\begin{minipage}[t]{7.45cm}\raggedright\sloppy\color{Ink}\fontsize{12.2}{16.1}\selectfont \par\vspace{.25cm}\lstinputlisting[language=R]{estratti/M13/durata.R}\end{minipage}\end{lrbox}
\begin{lrbox}{\RecallBox}\begin{minipage}[t]{16.45cm}\raggedright\sloppy\color{Ink}\fontsize{10.1}{13.2}\selectfont La durata viene derivata entro ciascuno scenario prima di riassumerne la distribuzione.\end{minipage}\end{lrbox}
\MeasureBox{27-TitleBox}{\TitleBox}
\MeasureBox{27-BodyBox}{\BodyBox}
\MeasureBox{27-SideBox}{\SideBox}
\MeasureBox{27-RecallBox}{\RecallBox}
\begin{frame}[plain]\hypertarget{slide27}{}
\bookmark[level=1,page=27]{Distribuzione degli scenari di durata}
\begin{tikzpicture}[remember picture,overlay]\node[anchor=north east,inner sep=0] at ([xshift=19.05cm,yshift=-0.16cm]current page.north west){\color{Classical}\fontsize{8}{10}\selectfont\hyperlink{slide2}{Indice}};
\node[anchor=north west,inner sep=0] at ([xshift=0.95cm,yshift=-0.62cm]current page.north west){\usebox{\TitleBox}};
\draw[Classical,line width=1.3pt] ([xshift=.95cm,yshift=-1.68cm]current page.north west)--([xshift=2.25cm,yshift=-1.68cm]current page.north west);\node[anchor=north west,inner sep=0] at ([xshift=0.95cm,yshift=-2.12cm]current page.north west){\usebox{\BodyBox}};
\node[anchor=north west,inner sep=0] at ([xshift=11.6cm,yshift=-2.1cm]current page.north west){\color{Quantum}\bfseries\fontsize{11.5}{14}\selectfont Implementazione in R};
\node[anchor=north west,inner sep=0] at ([xshift=11.6cm,yshift=-2.75cm]current page.north west){\usebox{\SideBox}};
\draw[Rule,line width=.5pt] ([xshift=.95cm,yshift=-10.43cm]current page.north west)--([xshift=19.05cm,yshift=-10.43cm]current page.north west);\node[anchor=north west,inner sep=0] at ([xshift=0.95cm,yshift=-10.66cm]current page.north west){\color{Classical}\bfseries\fontsize{9.4}{12}\selectfont Richiami};
\node[anchor=north west,inner sep=0] at ([xshift=2.6cm,yshift=-10.64cm]current page.north west){\usebox{\RecallBox}};
\draw[Rule,line width=.35pt] ([xshift=.95cm,yshift=-11.92cm]current page.north west)--([xshift=19.05cm,yshift=-11.92cm]current page.north west);\node[anchor=north west,inner sep=0] at ([xshift=0.95cm,yshift=-12.06cm]current page.north west){\color{Muted}\fontsize{7.7}{9.3}\selectfont M13 · Cronologie e archivi sedimentari};
\node[anchor=north east,inner sep=0] at ([xshift=17.65cm,yshift=-12.06cm]current page.north west){\color{Classical}\fontsize{7.3}{9.1}\selectfont };
\node[anchor=north east,inner sep=0] at ([xshift=19.05cm,yshift=-12.06cm]current page.north west){\color{Muted}\fontsize{7.7}{9.3}\selectfont 27/30};
\end{tikzpicture}\end{frame}
\begin{lrbox}{\TitleBox}\begin{minipage}[t]{18.2cm}\raggedright\sloppy\color{Ink}\fontsize{20.7}{24}\selectfont \bfseries\color{Navy}Confronto grafico delle durate ricostruite\end{minipage}\end{lrbox}
\begin{lrbox}{\BodyBox}\begin{minipage}[t]{10.25cm}\raggedright\sloppy\color{Ink}\fontsize{12.4}{16.4}\selectfont La distribuzione delle durate permette di distinguere il valore della cronologia di base dalla variabilità prodotta dalle ancore. La concentrazione degli scenari attorno a cinquecentotrenta anni riguarda l’intervallo selezionato e le deviazioni standard dichiarate. Intervalli con combinazioni diverse di ancore possono mostrare un’incertezza differente anche a parità di durata media. La lettura deve quindi collegare la forma dell’istogramma ai pesi dell’interpolazione, alla covarianza condivisa e alle fonti di variabilità che il modello ha lasciato fuori.\end{minipage}\end{lrbox}
\begin{lrbox}{\SideBox}\begin{minipage}[t]{7.45cm}\raggedright\sloppy\color{Ink}\fontsize{12.2}{16.1}\selectfont \includegraphics[width=7.45cm]{figure/M13/F11.pdf}\par\vspace{.16cm}{\fontsize{9.5}{12.6}\selectfont Durata fra 60 e 100 centimetri nei cinquecento scenari; la linea viola indica i 530 anni della cronologia di base.}\end{minipage}\end{lrbox}
\begin{lrbox}{\RecallBox}\begin{minipage}[t]{16.45cm}\raggedright\sloppy\color{Ink}\fontsize{10.1}{13.2}\selectfont L’incertezza di una durata dipende dalla collocazione delle ancore oltre che dall’ampiezza dell’intervallo.\end{minipage}\end{lrbox}
\MeasureBox{28-TitleBox}{\TitleBox}
\MeasureBox{28-BodyBox}{\BodyBox}
\MeasureBox{28-SideBox}{\SideBox}
\MeasureBox{28-RecallBox}{\RecallBox}
\begin{frame}[plain]\hypertarget{slide28}{}
\bookmark[level=1,page=28]{Confronto grafico delle durate ricostruite}
\begin{tikzpicture}[remember picture,overlay]\node[anchor=north east,inner sep=0] at ([xshift=19.05cm,yshift=-0.16cm]current page.north west){\color{Classical}\fontsize{8}{10}\selectfont\hyperlink{slide2}{Indice}};
\node[anchor=north west,inner sep=0] at ([xshift=0.95cm,yshift=-0.62cm]current page.north west){\usebox{\TitleBox}};
\draw[Classical,line width=1.3pt] ([xshift=.95cm,yshift=-1.68cm]current page.north west)--([xshift=2.25cm,yshift=-1.68cm]current page.north west);\node[anchor=north west,inner sep=0] at ([xshift=0.95cm,yshift=-2.12cm]current page.north west){\usebox{\BodyBox}};
\node[anchor=north west,inner sep=0] at ([xshift=11.6cm,yshift=-2.1cm]current page.north west){\color{Quantum}\bfseries\fontsize{11.5}{14}\selectfont Evidenza grafica};
\node[anchor=north west,inner sep=0] at ([xshift=11.6cm,yshift=-2.75cm]current page.north west){\usebox{\SideBox}};
\draw[Rule,line width=.5pt] ([xshift=.95cm,yshift=-10.43cm]current page.north west)--([xshift=19.05cm,yshift=-10.43cm]current page.north west);\node[anchor=north west,inner sep=0] at ([xshift=0.95cm,yshift=-10.66cm]current page.north west){\color{Classical}\bfseries\fontsize{9.4}{12}\selectfont Richiami};
\node[anchor=north west,inner sep=0] at ([xshift=2.6cm,yshift=-10.64cm]current page.north west){\usebox{\RecallBox}};
\draw[Rule,line width=.35pt] ([xshift=.95cm,yshift=-11.92cm]current page.north west)--([xshift=19.05cm,yshift=-11.92cm]current page.north west);\node[anchor=north west,inner sep=0] at ([xshift=0.95cm,yshift=-12.06cm]current page.north west){\color{Muted}\fontsize{7.7}{9.3}\selectfont M13 · Cronologie e archivi sedimentari};
\node[anchor=north east,inner sep=0] at ([xshift=17.65cm,yshift=-12.06cm]current page.north west){\color{Classical}\fontsize{7.3}{9.1}\selectfont };
\node[anchor=north east,inner sep=0] at ([xshift=19.05cm,yshift=-12.06cm]current page.north west){\color{Muted}\fontsize{7.7}{9.3}\selectfont 28/30};
\end{tikzpicture}\end{frame}
\begin{lrbox}{\TitleBox}\begin{minipage}[t]{18.2cm}\raggedright\sloppy\color{Ink}\fontsize{20.7}{24}\selectfont \bfseries\color{Navy}Collocazione del massimo osservato del proxy\end{minipage}\end{lrbox}
\begin{lrbox}{\BodyBox}\begin{minipage}[t]{10.25cm}\raggedright\sloppy\color{Ink}\fontsize{12.4}{16.4}\selectfont Il massimo numerico del proxy sintetico si trova a 160 centimetri, al bordo inferiore dell’archivio. La cronologia di base gli assegna 2200 anni cal BP, mentre lo scenario con hiatus lo colloca a 2600 anni. La differenza deriva da due assunzioni cronologiche applicate alla stessa osservazione e non documenta uno spostamento reale dell’evento. La posizione al bordo richiede inoltre prudenza nell’identificare un picco completo, poiché la serie potrebbe proseguire oltre l’intervallo campionato. Rumore, risoluzione e significato del proxy devono accompagnare l’interpretazione geologica.\par\vspace{.28cm}{\centering\fontsize{13}{17}\selectfont\(\displaystyle\begin{gathered}z_{\max}=\arg\max_{z_i}y_i,\qquad t_{\max}=t(z_{\max})\end{gathered}\)\par}\end{minipage}\end{lrbox}
\begin{lrbox}{\SideBox}\begin{minipage}[t]{7.45cm}\raggedright\sloppy\color{Ink}\fontsize{12.2}{16.1}\selectfont \input{tikz/M13/D09.tex}\end{minipage}\end{lrbox}
\begin{lrbox}{\RecallBox}\begin{minipage}[t]{16.45cm}\raggedright\sloppy\color{Ink}\fontsize{10.1}{13.2}\selectfont La collocazione temporale di un evento eredita le ipotesi della cronologia e l’incertezza del proxy.\end{minipage}\end{lrbox}
\MeasureBox{29-TitleBox}{\TitleBox}
\MeasureBox{29-BodyBox}{\BodyBox}
\MeasureBox{29-SideBox}{\SideBox}
\MeasureBox{29-RecallBox}{\RecallBox}
\begin{frame}[plain]\hypertarget{slide29}{}
\bookmark[level=1,page=29]{Collocazione del massimo osservato del proxy}
\begin{tikzpicture}[remember picture,overlay]\node[anchor=north east,inner sep=0] at ([xshift=19.05cm,yshift=-0.16cm]current page.north west){\color{Classical}\fontsize{8}{10}\selectfont\hyperlink{slide2}{Indice}};
\node[anchor=north west,inner sep=0] at ([xshift=0.95cm,yshift=-0.62cm]current page.north west){\usebox{\TitleBox}};
\draw[Classical,line width=1.3pt] ([xshift=.95cm,yshift=-1.68cm]current page.north west)--([xshift=2.25cm,yshift=-1.68cm]current page.north west);\node[anchor=north west,inner sep=0] at ([xshift=0.95cm,yshift=-2.12cm]current page.north west){\usebox{\BodyBox}};
\node[anchor=north west,inner sep=0] at ([xshift=11.6cm,yshift=-2.1cm]current page.north west){\color{Quantum}\bfseries\fontsize{11.5}{14}\selectfont Lettura del diagramma};
\node[anchor=north west,inner sep=0] at ([xshift=11.6cm,yshift=-2.75cm]current page.north west){\usebox{\SideBox}};
\draw[Rule,line width=.5pt] ([xshift=.95cm,yshift=-10.43cm]current page.north west)--([xshift=19.05cm,yshift=-10.43cm]current page.north west);\node[anchor=north west,inner sep=0] at ([xshift=0.95cm,yshift=-10.66cm]current page.north west){\color{Classical}\bfseries\fontsize{9.4}{12}\selectfont Richiami};
\node[anchor=north west,inner sep=0] at ([xshift=2.6cm,yshift=-10.64cm]current page.north west){\usebox{\RecallBox}};
\draw[Rule,line width=.35pt] ([xshift=.95cm,yshift=-11.92cm]current page.north west)--([xshift=19.05cm,yshift=-11.92cm]current page.north west);\node[anchor=north west,inner sep=0] at ([xshift=0.95cm,yshift=-12.06cm]current page.north west){\color{Muted}\fontsize{7.7}{9.3}\selectfont M13 · Cronologie e archivi sedimentari};
\node[anchor=north east,inner sep=0] at ([xshift=17.65cm,yshift=-12.06cm]current page.north west){\color{Classical}\fontsize{7.3}{9.1}\selectfont };
\node[anchor=north east,inner sep=0] at ([xshift=19.05cm,yshift=-12.06cm]current page.north west){\color{Muted}\fontsize{7.7}{9.3}\selectfont 29/30};
\end{tikzpicture}\end{frame}
\begin{lrbox}{\TitleBox}\begin{minipage}[t]{18.2cm}\raggedright\sloppy\color{Ink}\fontsize{20.7}{24}\selectfont \bfseries\color{Navy}Metadati cronologici e percorso di laboratorio\end{minipage}\end{lrbox}
\begin{lrbox}{\BodyBox}\begin{minipage}[t]{10.25cm}\raggedright\sloppy\color{Ink}\fontsize{12.4}{16.4}\selectfont La consegna conserva ancore originali, unità di profondità, convenzione BP, metodo di interpolazione, dominio, scenari di incertezza e decisioni di esclusione. I laboratori N01–N07 sviluppano controlli e tassi; N08–N12 propagano gli scenari e trasformano il proxy; N13–N16 esaminano discontinuità e sensibilità; N17–N20 completano pendenze, durate ed esportazione. Per lo studio autonomo si raddoppino le deviazioni standard e si confrontino numero di scenari accettati, età a cento centimetri e durata dell’intervallo sessanta–cento centimetri, motivando le differenze osservate.\end{minipage}\end{lrbox}
\begin{lrbox}{\SideBox}\begin{minipage}[t]{7.45cm}\raggedright\sloppy\color{Ink}\fontsize{12.2}{16.1}\selectfont \input{tikz/M13/D10.tex}\end{minipage}\end{lrbox}
\begin{lrbox}{\RecallBox}\begin{minipage}[t]{16.45cm}\raggedright\sloppy\color{Ink}\fontsize{10.1}{13.2}\selectfont La cronologia deve essere riproducibile come procedimento e interpretabile come insieme di ipotesi.\end{minipage}\end{lrbox}
\MeasureBox{30-TitleBox}{\TitleBox}
\MeasureBox{30-BodyBox}{\BodyBox}
\MeasureBox{30-SideBox}{\SideBox}
\MeasureBox{30-RecallBox}{\RecallBox}
\begin{frame}[plain]\hypertarget{slide30}{}
\bookmark[level=0,page=30]{Consegna e studio autonomo}
\bookmark[level=1,page=30]{Metadati cronologici e percorso di laboratorio}
\begin{tikzpicture}[remember picture,overlay]\node[anchor=north east,inner sep=0] at ([xshift=19.05cm,yshift=-0.16cm]current page.north west){\color{Classical}\fontsize{8}{10}\selectfont\hyperlink{slide2}{Indice}};
\node[anchor=north west,inner sep=0] at ([xshift=0.95cm,yshift=-0.62cm]current page.north west){\usebox{\TitleBox}};
\draw[Classical,line width=1.3pt] ([xshift=.95cm,yshift=-1.68cm]current page.north west)--([xshift=2.25cm,yshift=-1.68cm]current page.north west);\node[anchor=north west,inner sep=0] at ([xshift=0.95cm,yshift=-2.12cm]current page.north west){\usebox{\BodyBox}};
\node[anchor=north west,inner sep=0] at ([xshift=11.6cm,yshift=-2.1cm]current page.north west){\color{Quantum}\bfseries\fontsize{11.5}{14}\selectfont Lettura del diagramma};
\node[anchor=north west,inner sep=0] at ([xshift=11.6cm,yshift=-2.75cm]current page.north west){\usebox{\SideBox}};
\draw[Rule,line width=.5pt] ([xshift=.95cm,yshift=-10.43cm]current page.north west)--([xshift=19.05cm,yshift=-10.43cm]current page.north west);\node[anchor=north west,inner sep=0] at ([xshift=0.95cm,yshift=-10.66cm]current page.north west){\color{Classical}\bfseries\fontsize{9.4}{12}\selectfont Richiami};
\node[anchor=north west,inner sep=0] at ([xshift=2.6cm,yshift=-10.64cm]current page.north west){\usebox{\RecallBox}};
\draw[Rule,line width=.35pt] ([xshift=.95cm,yshift=-11.92cm]current page.north west)--([xshift=19.05cm,yshift=-11.92cm]current page.north west);\node[anchor=north west,inner sep=0] at ([xshift=0.95cm,yshift=-12.06cm]current page.north west){\color{Muted}\fontsize{7.7}{9.3}\selectfont M13 · Cronologie e archivi sedimentari};
\node[anchor=north east,inner sep=0] at ([xshift=17.65cm,yshift=-12.06cm]current page.north west){\color{Classical}\fontsize{7.3}{9.1}\selectfont };
\node[anchor=north east,inner sep=0] at ([xshift=19.05cm,yshift=-12.06cm]current page.north west){\color{Muted}\fontsize{7.7}{9.3}\selectfont 30/30};
\end{tikzpicture}\end{frame}
\end{document}